%% Fix harmless fixltx2e warning from elsarticle
\RequirePackage{silence}
\WarningFilter*{fixltx2e}{fixltx2e is not required}

\documentclass[5p,times]{elsarticle}
\raggedbottom

% =========================================================
% PACKAGES
% =========================================================

\usepackage{amsmath}
\usepackage{amssymb}
\usepackage{booktabs}
\usepackage{multirow}
\usepackage{graphicx}
\usepackage{grffile}
\usepackage{float}
\usepackage{tabularx}
\usepackage{array}
\usepackage{caption}
\usepackage{subcaption}
\usepackage{setspace}
\usepackage{enumitem}
\usepackage{url}
\usepackage{makecell}
\usepackage{xcolor}
\usepackage{placeins}
\usepackage{dblfloatfix}
\usepackage{pifont}
\usepackage[strings]{underscore}
\newcommand{\cmark}{\ding{51}}
\newcommand{\xmark}{\ding{55}}

\usepackage{hyperref}
\hypersetup{
	colorlinks=true,
	linkcolor=blue,
	citecolor=blue,
	urlcolor=blue
}

% microtype improves character protrusion/kerning and reduces the
% number and severity of overfull/underfull hboxes without altering
% the visible layout, wording, or float placement of the document.
\usepackage{microtype}

\allowdisplaybreaks
\emergencystretch=2em

% =========================================================
% SAFE FIGURE INCLUSION
% Prevents "Emergency stop" when a referenced image file
% has not yet been uploaded to the project. If the file
% exists, it displays normally. If not, a placeholder box
% is shown instead, and compilation continues.
% =========================================================
\newcommand{\safefig}[2]{%
	\IfFileExists{#1}{%
		\includegraphics[width=\linewidth]{#1}%
	}{%
		\fbox{\parbox[c][6cm][c]{0.9\linewidth}{\centering
				\textbf{Missing figure file:}\\ \texttt{#1}\\[4pt]
				\small (#2)}}%
	}%
}

% =========================================================
% PAGE & FLOAT TUNING
% =========================================================
\renewcommand{\topfraction}{0.9}
\renewcommand{\bottomfraction}{0.85}
\renewcommand{\textfraction}{0.05}
\renewcommand{\floatpagefraction}{0.9}

% ---------------------------------------------------------
% Consistent caption formatting (figures AND tables share the
% same font family/size, per formatting-consistency requirement)
% ---------------------------------------------------------
\captionsetup{font={small,rm},labelfont={small,bf,rm},labelsep=period}
\captionsetup[figure]{position=below,skip=4pt,font={small,rm},labelfont={small,bf,rm}}
\captionsetup[table]{position=above,skip=4pt,font={small,rm},labelfont={small,bf,rm}}

\setlength{\textfloatsep}{10pt plus 2pt minus 2pt}
\setlength{\floatsep}{8pt plus 2pt minus 2pt}
\setlength{\intextsep}{8pt plus 2pt minus 2pt}
\setlength{\abovecaptionskip}{4pt}
\setlength{\belowcaptionskip}{2pt}

% Heading spacing before/after \section, \subsection, and
% \subsubsection is governed uniformly by the document class's
% own definitions, so all headings at the same level already
% receive identical spacing throughout the document.

\newcolumntype{C}[1]{>{\centering\arraybackslash}p{#1}}
\newcolumntype{L}[1]{>{\raggedright\arraybackslash}p{#1}}

% =========================================================
% CUSTOM COLOURS
% =========================================================
\definecolor{barblue}{RGB}{31,119,180}
\definecolor{bargreen}{RGB}{44,160,44}
\definecolor{baryellow}{RGB}{255,187,0}
\definecolor{barred}{RGB}{214,39,40}
\definecolor{esncolor}{RGB}{72,176,230}
\definecolor{transcolor}{RGB}{255,140,0}
\definecolor{classcolor}{RGB}{80,200,120}
\definecolor{arrowgray}{RGB}{120,120,120}
\definecolor{darkblue}{RGB}{0,51,102}
\definecolor{medblue}{RGB}{31,119,180}
\definecolor{lightblue}{RGB}{174,214,241}
\definecolor{darkorange}{RGB}{200,100,0}
\definecolor{darkgreen}{RGB}{0,100,60}
\definecolor{darkred}{RGB}{150,30,30}
\definecolor{midgray}{RGB}{100,100,100}

% =========================================================
\begin{document}
	% =========================================================
	
	\begin{frontmatter}
		
		\title{A Reservoir Network Based Approach to Intrusion Detection
			Using Hybrid Transformer-ESN Architecture on CIC-IDS2017}
		
		\author[1]{L.~K.~Suresh Kumar}
		\author[1]{Jay Shreeram Yeraballi}
		\author[2]{Ravi Uyyala}
		\author[2]{Padmavathi Vurubindi}
		\author[3]{Ashok Kumar Das}
		
		\affiliation[1]{organization={Department of Computer Science and Engineering,
				University College of Engineering, Osmania University},
			city={Hyderabad}, country={India}}
		
		\affiliation[2]{organization={Department of Computer Science and Engineering,
				Chaitanya Bharathi Institute of Technology (CBIT)},
			postcode={500075}, city={Hyderabad}, country={India}}
		
		\affiliation[3]{organization={Center for Security, Theory and Algorithmic Research,
				International Institute of Information Technology},
			postcode={500032}, city={Hyderabad}, country={India}}
		
		\begin{abstract}
			
			Network intrusion detection under severe class imbalance remains challenging, as high aggregate accuracy can obscure poor recognition of minority classes. This paper proposes a hybrid Transformer--Echo State Network (Transformer-ESN) architecture for multi-class intrusion detection on CIC-IDS2017. The model combines a fixed, randomly initialised ESN reservoir with a multi-head self-attention Transformer encoder. The reservoir maps 76-dimensional preprocessed flow features into a 256-dimensional nonlinear state space without training recurrent weights, while the Transformer captures contextual dependencies among the resulting states. The trainable portion of the model is limited to the projection layer, Transformer encoder, and classification head, resulting in approximately 508K trainable parameters.
			
			Experiments using a stratified 80:20 train-test split show that Transformer-ESN achieves 99.38\% accuracy, 85.34\% macro-precision, 82.79\% macro-recall, and 73.55\% macro-F1. Although a reproduced Random Forest baseline achieves higher overall and macro-F1 performance on this tabular dataset, per-class analysis shows that the lower macro-F1 of Transformer-ESN is mainly due to four ultra-rare attack classes with very limited test support. Across the ten well-represented classes, the proposed model achieves F1-scores of at least 0.94 on nine.
			
			These findings suggest that reservoir-to-attention coupling is a stable and parameter-efficient design for intrusion detection. They also indicate that CICFlowMeter-based flow features are better suited to classical tabular learners than to sequential models, implying that the proposed architecture may be more effective for packet-level or temporally ordered traffic streams.
		\end{abstract}
		
		\begin{keyword}
			Intrusion Detection System \sep
			Echo State Network \sep
			Reservoir Computing \sep
			Transformer \sep
			Network Traffic Classification \sep
			CIC-IDS2017 \sep
			Deep Learning \sep
			Class Imbalance \sep
			Cybersecurity
		\end{keyword}
		
	\end{frontmatter}
	
	
	% =========================================================
	% 1. INTRODUCTION
	% =========================================================
	\section{Introduction}
	\label{sec:intro}
	
	Network security monitoring is now a core operational requirement as the
	attack surface of modern infrastructure grows in scale and sophistication.
	Organisations migrating services to cloud and software-defined networks
	face adversarial traffic that has outpaced manual analysis in volume and
	diversity~\cite{buczak2016survey,khraisat2019survey}. Network Intrusion
	Detection Systems (NIDS) must therefore classify high-throughput flows in
	real time and reliably detect attack types that appear infrequently in
	training data.
	The concept of intrusion detection was first formalized by Denning, who proposed one of the earliest intrusion detection models and laid the foundation for modern IDS research \cite{denning1987}.
	
	Deep learning has advanced automatic feature extraction for intrusion
	detection, reducing reliance on hand-engineered
	signatures~\cite{vinayakumar2019deep,yin2017deep}. Recurrent architectures
	capture temporal dependencies in traffic sequences~\cite{grurnn2018};
	convolutional models extract local feature
	correlations~\cite{cnnlstm2024,cnnbigru2022}; and Transformer encoders
	model long-range feature dependencies without
	recurrence~\cite{rtids2022,flowformer2024,long2024transformer}. The
	dominant benchmark across these paradigms is CIC-IDS2017, covering
	fourteen traffic classes under a controlled topology with realistic
	background traffic~\cite{cicids2017}.Deep neural network models combined with feature selection techniques have also been successfully applied to network intrusion detection using benchmark datasets such as NSL-KDD \cite{mlp2021}.
	
	Among traditional machine learning methods, Random Forest remains one of the most widely adopted classifiers because of its robustness, ensemble learning strategy, and strong generalization capability \cite{breiman2001}.
	IDS research has converged on an underexamined consensus: high aggregate
	accuracy on this benchmark is treated as evidence of practical utility.
	This is only partially justified. CIC-IDS2017 is severely
	imbalanced---benign traffic exceeds 80\% of flows, and several attack
	categories number fewer than fifty samples~\cite{cicids2017,ring2019survey}---so
	a classifier reporting 99\% accuracy may be near-chance on the minority
	categories that represent the most consequential threats: infiltration,
	targeted web exploits, and credential brute-force. Most published
	evaluations report aggregate accuracy alone, without per-class
	breakdowns, making it structurally impossible to verify whether any
	architecture detects rare attack
	traffic~\cite{thakkar2020,buczak2016survey}.
	
	This gap has direct operational consequences. A false negative on an
	Infiltration or SQL Injection event risks undetected data exfiltration or
	system compromise; misclassifying a benign flow merely triggers a
	recoverable false alarm. Aggregate accuracy conflates these asymmetric
	error costs into a single scalar dominated by the majority class, making
	macro-averaged F1 and per-class precision-recall analysis
	necessary---not merely preferable---for assessing deployability. This
	paper treats per-class analysis as a primary evaluation requirement and
	presents an honest characterisation of a novel hybrid architecture's
	behaviour across the full class distribution of CIC-IDS2017, compared
	against a strong classical baseline under identical conditions.
	
	\subsection{Literature Review}
	\label{sec:litreview}
	
	IDS research has evolved along four principal paradigms: classical
	machine learning, deep recurrent/convolutional networks, reservoir
	computing, and attention-based architectures. Table~\ref{tab:lit_comparison}
	summarises representative work across all four.
	
	\subsubsection{Classical Machine Learning}
	\label{subsec:classical}
	
	Classical approaches remain competitive despite their simplicity. Tree
	ensembles treat each flow independently yet routinely achieve
	near-perfect accuracy on CIC-IDS2017: 99.50\% with SMOTE-balanced Random
	Forest~\cite{kurniabudi2020}, 99.60\% with XGBoost~\cite{xgb2023},Gradient boosting methods such as XGBoost have also demonstrated excellent performance for intrusion detection owing to their scalability and handling of complex feature interactions \cite{chen2016} and 99.67\%
	with hybrid SVM-ELM~\cite{svmelm2018}. These results reflect the
	CICFlowMeter feature set~\cite{cicids2017}, whose pre-computed per-flow
	statistics already encode strong discriminative signal in tabular
	form---well suited to tree ensembles.
	
	This ceiling is structural: when attack behaviour spans multiple flows,
	or detection depends on temporal ordering that per-flow aggregation
	discards, flow-independent classifiers fail in ways further feature
	engineering cannot fix~\cite{buczak2016survey}. As elsewhere in this
	literature, evaluation is confined to aggregate accuracy, with no
	per-class breakdown indicating genuine minority-class detection versus
	accurate majority-class classification.
	
	\subsubsection{Deep Recurrent and Convolutional Networks}
	\label{subsec:deeplearning}
	
	Recurrent architectures assume packet and flow ordering encodes
	information absent from any single observation. Memory-equipped models
	outperform shallow classifiers on temporally structured
	traffic~\cite{vinayakumar2019deep,yin2017deep}, and gated GRUs
	substantially outperform vanilla RNNs on CIC-IDS2017 (99.00\% vs.\
	89.60\%)~\cite{grurnn2018}, attributed to improved gradient stability.
	
	Hybrid convolutional-recurrent designs separate local feature extraction
	from sequential modelling: CNN-BiGRU reaches
	99.70\%~\cite{cnnbigru2022} and CNN-LSTM reaches
	96.21\%~\cite{cnnlstm2024}, though both share the field's evaluation gap
	of aggregate-only accuracy. Convolutional networks applied directly to
	raw traffic also extract local patterns
	effectively~\cite{cnn2019}, though recurrent components remain
	necessary when cross-flow temporal ordering matters.
	
	\subsubsection{Reservoir Computing}
	\label{subsec:reservoir}
	
	Reservoir Computing (RC) offers gradient-free temporal feature extraction
	as an alternative to trained recurrent networks. The Echo State Network
	(ESN)~\cite{jaeger2001} uses a fixed, randomly initialised recurrent
	reservoir to map input sequences into a high-dimensional state space;
	provided the \emph{echo state property} holds---initial-condition
	influence decays asymptotically---a simple trained readout suffices for
	complex sequence classification~\cite{lukosevicius2012}, with the
	reservoir's spectral radius governing the stability-memory trade-off.
	Because reservoir weights are frozen, training reduces to fitting a
	single readout, avoiding vanishing gradients and benefiting streaming or
	resource-constrained settings.
	
	Prior ESN applications to intrusion detection share a limitation:
	standalone ESN classifiers achieve competitive accuracy on volumetric
	attacks, but their linear readouts cannot separate non-linearly
	distributed minority classes~\cite{esn_dos2021,mlesn2020}. Stacking
	reservoirs (DeepESN) improves representation capacity (97.60\% on
	CIC-IDS2017)~\cite{gallicchio2017}, and replacing the linear readout
	with a convolutional encoder plus hyperparameter optimisation raises
	readout expressiveness further: a CNN-ESN with chaotic walrus
	optimisation reports 99.80\% accuracy, though this figure is reported
	on NSL-KDD, WSN-DS, and IoT-23 rather than CIC-IDS2017 and is therefore
	not directly comparable to the CIC-IDS2017 results discussed
	elsewhere in this review~\cite{cesn2024}. Taken together, these results
	suggest that readout expressiveness, not the reservoir itself, is the
	binding constraint. This motivates the present work's use of the
	reservoir as a gradient-free encoder paired with a Transformer readout.
	
	\subsubsection{Attention-Based Architectures}
	\label{subsec:attention}
	
	The Transformer~\cite{vaswani2017attention} replaces sequential
	recurrence with parallel self-attention, capturing arbitrary-range
	dependencies in a single forward pass while multi-head attention attends
	jointly to different input aspects.
	
	Applications to IDS have been competitive: a standalone Transformer on
	flow-level features reached 99.17\%~\cite{rtids2022}; multi-head
	attention combined with BiLSTM captured long-range dependencies that
	recurrent gating misses~\cite{bilstm_attn2023}; FlowTransformer
	incorporated domain-specific inductive biases for flow-based
	IDS~\cite{flowformer2024}; and TRBMA traded aggregate accuracy (98.66\%)
	for higher macro-recall via bidirectional temporal attention, explicitly
	prioritising minority-class detection~\cite{trbma2025}. A recent survey
	maps the broader Transformer and LLM landscape in intrusion
	detection~\cite{kheddar2024transformers}. A shared limitation across
	this paradigm is the absence of gradient-free temporal pre-processing:
	Transformers operate on fixed-length vectors and capture inter-feature
	but not cross-flow temporal structure unless an explicit sequential
	representation is built first.
	
	\subsubsection{Research Gap}
	\label{subsec:gap}
	
	The literature above establishes two gaps. \textbf{Architectural gap:}
	no prior work pairs a fixed ESN reservoir with a Transformer
	encoder---ESN reservoirs have been evaluated only with linear or shallow
	readouts, and Transformers applied without gradient-free temporal
	front-ends. \textbf{Evaluation gap:} published IDS evaluations report
	aggregate accuracy without per-class breakdown under CIC-IDS2017's
	severe imbalance, making it impossible to distinguish genuine
	minority-class detection from majority-class-driven scores. This work
	reports per-class precision, recall, and F1 for all fourteen categories
	alongside macro-averaged metrics. Independent 2025--2026 work confirms
	this gap remains unresolved: ensemble methods targeting minority-class
	detection report persistent difficulty on ultra-rare
	categories~\cite{oh2025mchensemble}, and class-imbalance mitigation
	frameworks report similar limitations on comparable
	benchmarks~\cite{hambali2026breaking}, though the latter used NSL-KDD
	rather than CIC-IDS2017.
	
	\begin{table*}[!htbp]
		\centering
		\caption{Representative work on CIC-IDS2017 across the four paradigms
			reviewed. ``Per-class'' indicates whether per-class breakdown was
			reported. ``Temporal'' indicates whether the architecture explicitly
			models cross-flow or within-sequence temporal structure.
			$\dagger$~Macro-recall specifically prioritised by design.}
		\label{tab:lit_comparison}
		\resizebox{\textwidth}{!}{%
			\begin{tabular}{L{3.5cm} L{3.5cm} C{2.2cm} C{2.2cm} C{1.8cm} C{1.8cm}}
				\toprule
				\textbf{Reference} & \textbf{Architecture} & \textbf{Accuracy (\%)}
				& \textbf{Macro-F1} & \textbf{Per-class} & \textbf{Temporal} \\
				\midrule
				\multicolumn{6}{l}{\textit{Classical ML}} \\
				Pranto et al.~\cite{kurniabudi2020}       & Random Forest + SMOTE         & 99.50 & ---  & No  & No  \\
				Azam et al.~\cite{xgb2023}        & XGBoost                       & 99.60 & ---  & No  & No  \\
				Ahmad et al.~\cite{svmelm2018}    & SVM-ELM                       & 99.67 & ---  & No  & No  \\
				\midrule
				\multicolumn{6}{l}{\textit{Deep Recurrent / Convolutional}} \\
				Tang et al.~\cite{grurnn2018}     & GRU-RNN                       & 99.00 & ---  & No  & Yes \\
				Wang et al.~\cite{cnnbigru2022}   & CNN-BiGRU                     & 99.70 & ---  & No  & Yes \\
				Kaur et al.~\cite{cnnlstm2024}    & CNN-LSTM                      & 96.21 & ---  & No  & Yes \\
				\midrule
				\multicolumn{6}{l}{\textit{Reservoir Computing}} \\
				Gallicchio et al.~\cite{gallicchio2017} & DeepESN               & 97.60 & ---  & No  & Yes \\
				\midrule
				\multicolumn{6}{l}{\textit{Attention-Based}} \\
				Wu et al.~\cite{rtids2022}        & Transformer                   & 99.17 & ---  & No  & No  \\
				Zhang et al.~\cite{bilstm_attn2023} & BiLSTM + Attention          & 99.40 & ---  & No  & Yes \\
				Li et al.~\cite{trbma2025}        & TRBMA                         & 98.66 & High$\dagger$ & No & Yes \\
				\midrule
				\multicolumn{6}{l}{\textit{Proposed}} \\
				\textbf{This work}                & \textbf{Transformer-ESN}      & \textbf{99.38} & \textbf{Reported} & \textbf{Yes} & \textbf{Yes} \\
				\bottomrule
		\end{tabular}}
	\end{table*}
	
	Figure~\ref{fig:accuracy_comparison} ranks the reported accuracy of
	these representative classifiers to situate the proposed model's
	performance within this range visually.
	
	\begin{figure}[!htbp]
		\centering
		\safefig{Figure1_Accuracy_Ranking_Elsevier.png}{Accuracy ranking chart across four learning paradigms}
		\caption{Ranked comparison of reported accuracy for representative CIC-IDS2017 classifiers spanning classical machine learning, deep recurrent/convolutional networks, reservoir computing, and attention-based architectures. The proposed Transformer-ESN (99.38\%) is positioned within the 96.21\%--99.70\% range established by prior work. As protocol heterogeneity across studies precludes direct numerical comparison~\cite{thakkar2020}, the ranking is provided for contextual orientation rather than head-to-head benchmarking.}
		\label{fig:accuracy_comparison}
	\end{figure}
	
	As Figure~\ref{fig:accuracy_comparison} shows, the proposed model's
	accuracy sits inside the range already established by prior classifiers;
	Section~\ref{subsec:context} returns to this comparison in
	detail once the full results are presented.
	
	\subsection{The Proposed Architecture: Transformer-ESN}
	
	The architecture proposed here---\textbf{Transformer-ESN}---couples a
	fixed Echo State Network (ESN) reservoir with a multi-head self-attention
	Transformer encoder. In ESNs, the recurrent weight matrix is fixed at
	random initialisation and only the readout is
	trained~\cite{jaeger2001,lukosevicius2012}, eliminating vanishing and
	exploding gradients and reducing training
	complexity~\cite{gallicchio2017,scardapane2017}. Prior ESN-based IDS
	work has used the reservoir's nonlinear states directly as classifier
	inputs~\cite{esn_dos2021,mlesn2020,cesn2024}.
	
	Transformer-ESN instead treats the reservoir as a structured nonlinear
	projection layer: it maps 76-dimensional preprocessed CICFlowMeter
	features into a 256-dimensional nonlinear state space at zero gradient
	cost. The resulting states are passed as a token sequence to a
	multi-head self-attention encoder, which learns global contextual
	dependencies over the projected feature
	space~\cite{vaswani2017attention}, with a softmax head producing
	per-class probabilities over the fourteen categories. The trainable
	component---encoder weights and classification head---totals
	approximately 508K parameters. To our knowledge, this coupling has not
	been previously reported.
	
	\subsection{The Structural Mismatch Finding}
	
	Transformer-ESN achieves 99.38\% overall accuracy, competitive across
	all four surveyed paradigms. However, a Random Forest baseline under
	identical preprocessing and split achieves 99.80\% accuracy and 88.40\%
	macro-F1, versus 73.55\% macro-F1 for Transformer-ESN.
	
	The more informative finding concerns the \textit{source} of this gap.
	CICFlowMeter features are pre-aggregated statistics computed only after
	a flow terminates~\cite{cicids2017}, so each row is an independent
	observation rather than a step in a temporal sequence. Sequential
	architectures (ESNs, LSTMs, GRUs, and the hybrid proposed here) apply a
	sequential inductive bias to data that lacks the row-wise ordering that
	bias exploits. Random Forest makes no such assumption, treating each row
	as an independent point in 76-dimensional space---precisely what
	CICFlowMeter produces~\cite{kurniabudi2020,xgb2023}. This structural mismatch,
	not a difference in representational capacity, is the primary driver of
	the macro-F1 gap. The same sequential bias becomes an advantage where
	input genuinely consists of ordered sequences---packet-level streams,
	session logs, or raw byte sequences---where reservoir dynamics can
	exploit real temporal structure.
	
	The gap is amplified by extreme rarity in four categories: Web Attack --
	SQL Injection has four test samples and Infiltration has seven, so a
	single misclassification shifts per-class F1 by 25\% and 14\%
	respectively---these metrics are not statistically stable estimators of
	architectural capability at this support level. Across the ten
	well-supported classes, Transformer-ESN achieves F1 $\geq$ 0.94 on nine,
	confirming correct performance where sufficient training signal exists.
	
	\subsection{Scope and Contributions}
	
	This work makes the following contributions:
	
	\begin{itemize}[leftmargin=15pt,itemsep=3pt,topsep=4pt]
		
		\item \textbf{A novel hybrid architecture} coupling a fixed ESN
		reservoir with a Transformer encoder for multi-class network
		intrusion detection, implemented as a 508K-parameter trainable
		model with recurrent weight training eliminated from the feature
		encoding stage.
		
		\item \textbf{An empirically grounded per-class evaluation} on
		CIC-IDS2017 demonstrating competitive overall accuracy (99.38\%)
		alongside an explicit characterisation of failure modes on
		ultra-rare attack classes (test support $\leq$ 301 samples)---a
		level of transparency absent from most published IDS evaluations.
		
		\item \textbf{A direct comparison against a reproduced Random Forest
			baseline} under identical preprocessing, split, and evaluation
		protocols, establishing that the macro-F1 gap (73.55\% vs.\
		88.40\%) is attributable to ultra-rare class behaviour and
		structural mismatch, not systematic architectural inferiority on
		well-represented traffic categories.
		
		\item \textbf{A structured analysis of architectural mismatch}
		between sequential inductive biases and pre-aggregated tabular flow
		features, identifying where reservoir-attention hybrids are
		expected to add value over classical tabular methods, and where
		they do not.
		
		\item \textbf{A complete theoretical treatment} of the pipeline,
		covering ESN reservoir dynamics, the echo state property and
		spectral radius conditioning, positional encoding over reservoir
		tokens, scaled dot-product attention, and multi-head projection.
		
	\end{itemize}
	
	The remainder of this paper is organised as follows.
	Section~\ref{sec:methodology} describes the dataset, preprocessing
	pipeline, and complete architectural specification.
	Several studies have performed feature analysis on the CICIDS2017 dataset using Information Gain to improve anomaly detection performance \cite{rf2020}.
	Section~\ref{sec:results} presents experimental results with full
	per-class analysis. Section~\ref{sec:discussion} provides comparative
	analysis and a structured decomposition of the performance gap.
	Section~\ref{sec:conclusion} concludes and outlines future directions.
	
	% =========================================================
	% 3. METHODOLOGY
	% =========================================================
	\section{Methodology}
	\label{sec:methodology}
	
	\subsection{Dataset and Challenges}
	\label{subsec:dataset}
	
	This work uses the CIC-IDS2017 benchmark~\cite{cicids2017}, one of the
	most widely referenced datasets in network intrusion detection
	research. Produced by the Canadian Institute for Cybersecurity, it was
	collected over five consecutive days in a controlled testbed that
	closely replicates an operational organisational network, and covers
	fourteen traffic classes: one benign class and thirteen attack
	categories.Previous studies have also investigated the importance of individual CICIDS2017 features to identify the most discriminative attributes for attack detection \cite{reis2019}.Although CICIDS2017 is widely adopted, researchers have reported several quality issues and preprocessing challenges that should be carefully addressed before model training \cite{engelen2021}.
	The flow-based features used in CICIDS2017 originate from the CICFlowMeter framework proposed by Lashkari et al. \cite{lashkari2017}.
	
	\subsubsection{Dataset Characteristics}
	
	The thirteen attack categories include DDoS, DoS Hulk, DoS GoldenEye,
	DoS Slowloris, DoS Slowhttptest, Botnet ARES, PortScan, FTP-Patator,
	SSH-Patator, Web Attack Brute Force, Web Attack SQL Injection, Web
	Attack XSS, and Infiltration. Traffic was characterised using
	CICFlowMeter~\cite{cicids2017}, which produces 78 flow-level
	statistical attributes from packet lengths, inter-arrival times, flow
	durations, TCP flag counts, and subflow byte/packet statistics---the
	same properties exploited by classical ML methods~\cite{kurniabudi2020,xgb2023},
	while offering enough structural richness to benefit more expressive
	architectures. Key properties are summarised in
	Table~\ref{tab:dataset}.
	
	\begin{table}[!htbp]
		\centering
		\caption{Summary of CIC-IDS2017 dataset properties.}
		\label{tab:dataset}
		\renewcommand{\arraystretch}{0.9}
		\setlength{\tabcolsep}{5pt}
		\resizebox{\linewidth}{!}{%
			\begin{tabular}{ll}
				\toprule
				\textbf{Property}          & \textbf{Description} \\
				\midrule
				Dataset                    & CIC-IDS2017 \\
				Features                   & 78 (flow-level statistical) \\
				Classes                    & 14 (1 Benign + 13 Attack) \\
				Traffic Type               & Benign + Malicious \\
				Classification Task        & Multi-class \\
				Feature Extraction Tool    & CICFlowMeter \\
				Collection Period          & Monday to Friday (5 days) \\
				Environment                & Controlled realistic testbed \\
				Attack Categories          & DDoS, DoS (4 variants), Botnet, PortScan, \\
				& FTP-Patator, SSH-Patator, \\
				& Web Attacks (3 variants), Infiltration \\
				Evaluation Metrics         & Accuracy, Macro-Prec., Macro-Rec., Macro-F1 \\
				\bottomrule
		\end{tabular}}
	\end{table}
	
	\subsubsection{Class Distribution and Imbalance}
	
	Class imbalance is the dataset's most significant challenge.
	High-volume classes such as DoS Hulk and PortScan each contribute
	hundreds of thousands of samples, while minority classes such as
	Infiltration (36 samples) account for less than 0.01\% of the total.
	Table~\ref{tab:class_dist} gives the full distribution. This skew is
	considerably more severe than in comparable benchmarks such as NSL-KDD
	and UNSW-NB15, making CIC-IDS2017 a demanding test of minority-class
	detection.
	
	\begin{table}[!htbp]
		\centering
		\caption{Approximate class distribution in CIC-IDS2017. Minority
			classes ($<$0.1\%) represent the primary evaluation challenge.}
		\label{tab:class_dist}
		\renewcommand{\arraystretch}{0.88}
		\setlength{\tabcolsep}{5pt}
		\resizebox{\linewidth}{!}{%
			\begin{tabular}{L{3.8cm} C{2.0cm} C{2.0cm}}
				\toprule
				\textbf{Class}
				& \textbf{Approx.\ Samples}
				& \textbf{Relative Freq.\ (\%)} \\
				\midrule
				Benign                    & 2{,}271{,}320 & 80.30 \\
				DoS Hulk                  &   231{,}073   &  8.17 \\
				PortScan                  &   158{,}930   &  5.62 \\
				DDoS                      &   128{,}027   &  4.53 \\
				DoS GoldenEye             &    10{,}293   &  0.36 \\
				FTP-Patator               &     7{,}938   &  0.28 \\
				SSH-Patator               &     5{,}897   &  0.21 \\
				DoS Slowloris             &     5{,}796   &  0.21 \\
				DoS Slowhttptest          &     5{,}499   &  0.19 \\
				Botnet                    &     1{,}966   &  0.07 \\
				Web Attack -- Brute Force &     1{,}507   &  0.05 \\
				Web Attack -- XSS         &       652     &  0.02 \\
				Web Attack -- SQL Inject  &        21     & $<$0.01 \\
				Infiltration              &        36     & $<$0.01 \\
				\bottomrule
		\end{tabular}}
	\end{table}
	
	Because a classifier can reach deceptively high aggregate accuracy by
	ignoring rare classes, macro-averaged precision, recall, and F1-score
	are treated as essential evaluation criteria alongside overall
	accuracy~\cite{trbma2025}; missing even one minority-class threat such
	as Infiltration is a serious operational failure regardless of the
	overall accuracy figure. To address the imbalance, this study applies
	class-weighted loss and per-class reporting, and assesses all models
	under the same four-metric protocol---accuracy, macro-precision,
	macro-recall, macro-F1---for fair comparison with the
	literature~\cite{cnnbigru2022,rtids2022,trbma2025}.
	
	\subsubsection{Feature Space Challenges}
	
	Beyond class imbalance, the 78-dimensional CICFlowMeter feature space
	introduces two further preprocessing issues: zero-duration flows
	produce undefined (infinite) values in \textit{Flow Bytes/s} and
	\textit{Flow Packets/s}, and high inter-feature correlation can
	mislead gradient-based optimisers toward spurious associations. These
	are addressed in Section~\ref{subsec:preprocessing}.
	
	\subsection{Proposed Transformer-ESN Framework}
	\label{subsec:arch}
	
	The proposed framework processes traffic through three stages:
	(i)~preprocessing and normalisation, (ii)~nonlinear temporal feature
	expansion via an ESN reservoir, and (iii)~contextual refinement and
	classification via a Transformer encoder, as shown in
	Figure~\ref{fig:arch}.
	
	\subsubsection{System Architecture Overview}
	
	Figure~\ref{fig:arch} summarises the end-to-end pipeline described in
	this subsection. Raw flow records first pass through the
	preprocessing and normalisation stage detailed in
	Section~\ref{subsec:preprocessing}, producing the 76-dimensional input
	vector $\mathbf{u} \in \mathbb{R}^{76}$. This vector is then projected
	through the frozen ESN reservoir into a 256-dimensional nonlinear
	state, passed through the trainable Transformer encoder for contextual
	refinement, and finally reduced by the classification head to a
	14-class probability distribution. The frozen and trainable stages are
	shown separately in the figure to make explicit which parameters are
	updated during backpropagation and which remain fixed at
	initialisation.
	
	\begin{figure}[!htbp]
		\centering
		\safefig{Figure2_Architecture_Comparison.png}{Transformer-ESN architecture pipeline}
		\caption{End-to-end architecture of the proposed Transformer-ESN framework. The frozen ESN reservoir maps the preprocessed 76-dimensional input into a 256-dimensional nonlinear state, which is projected, positionally encoded, and refined by a multi-head self-attention Transformer encoder before the classification head produces per-class probabilities. Reservoir weights $W_{\mathrm{in}}$ and $W$ remain fixed throughout training; only the projection layer, Transformer encoder, and classification head are updated via backpropagation.}
		\label{fig:arch}
	\end{figure}
	
	The four evaluation metrics used throughout this study are further
	summarised in radar form in Figure~\ref{fig:radar_chart}, which
	visualises the asymmetry between accuracy and macro-F1 alongside the
	paired grouped comparison in Figure~\ref{fig:four_metric}.
	
	\begin{figure}[!htbp]
		\centering
		\safefig{Figure3_Grouped_DotPlot.png}{Four-metric grouped comparison plot}
		\caption{Grouped comparison of accuracy, macro-precision, macro-recall, and macro-F1 for the reproduced Random Forest baseline and the proposed Transformer-ESN on CIC-IDS2017. The accuracy gap between the two methods is modest (0.42 percentage points), whereas the macro-F1 gap (14.85 percentage points) is substantially larger, reflecting the decomposition discussed in Section~\ref{subsec:decompose}. Precision and recall values for Random Forest were not reported in the reproduced baseline and are estimated from comparable configurations in the literature~\cite{kurniabudi2020}.}
		\label{fig:four_metric}
	\end{figure}
	
	\begin{figure}[!htbp]
		\centering
		\safefig{Figure4_TransformerESN_RadarChart.png}{Transformer-ESN multi-metric radar chart}
		\caption{Radar-chart representation of the proposed Transformer-ESN across the four evaluation dimensions used throughout this study: accuracy, macro-precision, macro-recall, and macro-F1. The asymmetry between the outer (accuracy) and inner (macro-F1) vertices visually summarises the class-imbalance effect quantified in Table~\ref{tab:overall} and decomposed in Section~\ref{subsec:decompose}.}
		\label{fig:radar_chart}
	\end{figure}
	
	The outer-to-inner gap visible in Figure~\ref{fig:radar_chart} between
	the accuracy vertex and the macro-F1 vertex is the same class-imbalance
	effect quantified numerically in Table~\ref{tab:overall} and decomposed
	in Section~\ref{subsec:decompose}.
	
	\subsubsection{Echo State Network: Theoretical Foundations}
	
	Echo State Networks (ESNs) belong to the Reservoir Computing
	family~\cite{jaeger2001,lukosevicius2012}. Unlike conventional
	recurrent networks trained via Backpropagation Through Time
	(BPTT)~\cite{grurnn2018}, an ESN uses a large randomly initialised
	recurrent reservoir whose weights are never updated; only the readout
	is trained, which eliminates vanishing/exploding gradients in the
	recurrent component entirely.
	
	\paragraph{Reservoir State Dynamics.}
	Let $\mathbf{u}(t) \in \mathbb{R}^{d_{\mathrm{in}}}$ be the input
	feature vector at time step $t$, with $d_{\mathrm{in}} = 76$ after
	preprocessing (Section~\ref{subsec:preprocessing}). The reservoir
	state is updated as:
	
	\begin{equation}
		\tilde{\mathbf{x}}(t) =
		f\!\left(W_{\mathrm{in}}\,\mathbf{u}(t) +
		W\,\mathbf{x}(t-1)\right)
		\label{eq:esn_update}
	\end{equation}
	
	\begin{equation}
		\mathbf{x}(t) =
		(1-\alpha)\,\mathbf{x}(t-1) + \alpha\,\tilde{\mathbf{x}}(t)
		\label{eq:esn_leaky}
	\end{equation}
	
	\noindent where $\mathbf{x}(t) \in \mathbb{R}^{N}$ is the reservoir
	state ($N = 256$), $W_{\mathrm{in}} \in \mathbb{R}^{N \times
		d_{\mathrm{in}}}$ and $W \in \mathbb{R}^{N \times N}$ are fixed input
	and recurrent weight matrices, $f(\cdot) = \tanh(\cdot)$, and $\alpha
	\in (0,1]$ is the leaking rate. Setting $\alpha = 1.0$, as done here,
	reduces Eq.~\eqref{eq:esn_leaky} to Eq.~\eqref{eq:esn_update},
	preserving the full instantaneous response without temporal
	smoothing.
	
	\paragraph{Echo State Property and Spectral Radius.}
	The Echo State Property (ESP) requires the effect of any initial state
	$\mathbf{x}(0)$ to decay asymptotically to zero, so reservoir outputs
	depend only on recent input history~\cite{jaeger2001}. The ESP holds
	when the spectral radius of $W$ satisfies $\rho(W) < 1$:
	
	\begin{equation}
		\rho(W) =
		\max\!\left\{|\lambda| : \lambda \text{ is an eigenvalue of }
		W\right\}
		\label{eq:spectral}
	\end{equation}
	
	$W$ is initialised as a sparse random matrix, then rescaled to a
	target spectral radius $r$:
	
	\begin{equation}
		W \leftarrow
		\frac{r}{\rho(W_{\mathrm{init}})} \cdot W_{\mathrm{init}}
		\label{eq:rescale}
	\end{equation}
	
	\noindent This work uses $r = 0.9$, satisfying the ESP while retaining
	sufficient reservoir memory for temporal patterns in network traffic,
	following the initialisation guidelines of Scardapane and
	Wang~\cite{scardapane2017}.
	
	\paragraph{Reservoir Capacity and Input Scaling.}
	Reservoir behaviour is governed primarily by size $N$, spectral radius
	$r$, and input scaling $\sigma$. An input scaling of $\sigma = 0.1$
	was chosen to keep a moderate nonlinear operating regime suited to
	CIC-IDS2017's statistical flow features. The reservoir expands the
	76-dimensional input into a 256-dimensional nonlinear state
	(${\approx}3.4\times$ expansion). Sparse connectivity (10\%) promotes
	unit diversity while preserving the lateral interactions needed for
	temporal integration~\cite{lukosevicius2012}.
	
	\subsubsection{Transformer Encoder: Theoretical Foundations}
	
	The Transformer encoder~\cite{vaswani2017attention} refines the
	sequence of ESN reservoir states through multi-head scaled
	dot-product self-attention, producing globally coherent contextual
	representations---the central architectural contribution relative to
	standard ESNs, which use a simple linear readout instead.
	
	\paragraph{Input Projection and Positional Encoding.}
	Each reservoir state $\mathbf{x}(t)$ is projected into the Transformer
	model dimension $d_{\mathrm{model}}$ via a learned linear map:
	
	\begin{equation}
		\mathbf{z}_t =
		W_{\mathrm{proj}}\,\mathbf{x}(t) + \mathbf{b}_{\mathrm{proj}},
		\quad \mathbf{z}_t \in \mathbb{R}^{d_{\mathrm{model}}}
		\label{eq:proj}
	\end{equation}
	
	Sinusoidal positional encodings~\cite{vaswani2017attention} are added
	to retain sequence order:
	
	\begin{align}
		\mathrm{PE}_{(t,\,2i)}   &=
		\sin\!\left(\frac{t}{10000^{2i/d_{\mathrm{model}}}}\right)
		\label{eq:pe_sin} \\
		\mathrm{PE}_{(t,\,2i+1)} &=
		\cos\!\left(\frac{t}{10000^{2i/d_{\mathrm{model}}}}\right)
		\label{eq:pe_cos}
	\end{align}
	
	Manocchio et al.~\cite{flowformer2024} showed that preserving
	positional information across flow-based token sequences measurably
	improves detection accuracy. The final encoded input is
	$\mathbf{h}_t = \mathbf{z}_t + \mathrm{PE}_t$.
	
	\paragraph{Scaled Dot-Product Attention.}
	For a single head, query, key, and value matrices are obtained from
	$H = [\mathbf{h}_1, \ldots, \mathbf{h}_T]$ via learned projections:
	
	\begin{equation}
		Q = HW^Q, \quad K = HW^K, \quad V = HW^V
		\label{eq:qkv}
	\end{equation}
	
	\begin{equation}
		\mathrm{Attention}(Q,K,V) =
		\mathrm{Softmax}\!\!\left(\frac{QK^{\top}}{\sqrt{d_k}}\right)V
		\label{eq:attention}
	\end{equation}
	
	where $d_k = d_{\mathrm{model}} / H_{\mathrm{heads}}$. The
	$1/\sqrt{d_k}$ scaling prevents dot-products from pushing the softmax
	into near-zero gradient regions~\cite{vaswani2017attention}.
	
	\paragraph{Multi-Head Self-Attention.}
	
	\begin{equation}
		\mathrm{head}_h =
		\mathrm{Attention}(HW^Q_h,\; HW^K_h,\; HW^V_h)
		\label{eq:head}
	\end{equation}
	
	\begin{equation}
		\mathrm{MHA}(H) =
		\mathrm{Concat}(\mathrm{head}_1,\ldots,\mathrm{head}_H)\,W^O
		\label{eq:mha}
	\end{equation}
	
	Four heads ($H = 4$) are used with $d_{\mathrm{model}} = 128$, giving
	$d_k = 32$ per head. Parallel heads let the model attend to
	complementary feature relationships in distinct subspaces, which is
	useful for distinguishing structurally similar attack families such
	as the four DoS variants in CIC-IDS2017.
	
	\paragraph{Position-wise Feed-Forward Network.}
	
	\begin{equation}
		\mathrm{FFN}(\mathbf{x}) =
		\max(0,\; \mathbf{x}W_1 + \mathbf{b}_1)\,W_2 + \mathbf{b}_2
		\label{eq:ffn}
	\end{equation}
	
	The inner projection uses $d_{\mathrm{ff}} = 256$ before projecting
	back to $d_{\mathrm{model}} = 128$, applied independently at each
	token position.
	
	\paragraph{Layer Normalisation and Residual Connections.}
	
	\begin{equation}
		\mathbf{h}'_t =
		\mathrm{LayerNorm}\!\left(\mathbf{h}_t +
		\mathrm{Sublayer}(\mathbf{h}_t)\right)
		\label{eq:layernorm}
	\end{equation}
	
	where $\mathrm{Sublayer}(\cdot)$ is either the MHA block
	(Eq.~\eqref{eq:mha}) or the FFN (Eq.~\eqref{eq:ffn}). Two encoder
	layers ($L = 2$) balance representational capacity against training
	complexity.
	
	\subsubsection{Hybrid Coupling Mechanism}
	
	The ESN and Transformer form a unified pipeline. Each input flow
	vector $\mathbf{u}(t)$ is transformed into a 256-dimensional
	reservoir state via Eqs.~\eqref{eq:esn_update}--\eqref{eq:esn_leaky},
	projected into the Transformer input space (Eq.~\eqref{eq:proj}), and
	enriched with positional encodings
	(Eqs.~\eqref{eq:pe_sin}--\eqref{eq:pe_cos}). Two layers of multi-head
	self-attention and FFN sublayers, each with residual connections and
	layer normalisation, refine the resulting tokens. Mean pooling
	followed by a fully connected softmax head yields the final per-class
	probability vector:
	
	\begin{equation}
		\hat{y} = \mathrm{Softmax}\!\!\left(
		f_{\mathrm{cls}}\!\!\left(
		\mathrm{MHA}\!\left(
		W_{\mathrm{proj}} \cdot f_{\mathrm{ESN}}(\mathbf{u})
		\right)
		\right)
		\right)
		\label{eq:endtoend}
	\end{equation}
	
	where $f_{\mathrm{ESN}}(\mathbf{u}) = \tanh(W_{\mathrm{in}}\,
	\mathbf{u} + W\,\mathbf{x}_{\mathrm{prev}})$. $W_{\mathrm{in}}$ and
	$W$ are frozen after random initialisation; only $W_{\mathrm{proj}}$,
	the Transformer encoder weights, and the classification head are
	updated during training.
	
	\subsubsection{Training Procedure}
	
	Figure~\ref{fig:training_flow} shows the training workflow. Reservoir
	weights stay frozen throughout; backpropagation passes only through
	the Transformer encoder and classification head, eliminating
	recurrent gradient computation and enabling stable convergence at
	lower training cost than fully trainable recurrent models.
	
	\subsubsection{Architectural Comparison}
	
	Table~\ref{tab:arch_compare} compares Transformer-ESN against
	representative architectures. It is the only evaluated model that
	both eliminates trainable recurrence via the fixed reservoir and
	incorporates multi-head self-attention, achieving a Low--Medium
	training cost not matched by any other attention-based architecture
	here.
	
	\begin{table}[!htbp]
		\centering
		\caption{Architectural comparison of the proposed Transformer-ESN
			with alternative paradigms on key design properties.}
		\label{tab:arch_compare}
		\renewcommand{\arraystretch}{0.9}
		\setlength{\tabcolsep}{4pt}
		\resizebox{\linewidth}{!}{%
			\begin{tabular}{L{2.6cm} C{1.4cm} C{1.3cm} C{1.4cm} C{1.3cm} C{1.5cm}}
				\toprule
				\textbf{Architecture}
				& \textbf{Trainable Recur.}
				& \textbf{Self-Attn.}
				& \textbf{Fixed Res.}
				& \textbf{Quad.\ Attn.}
				& \textbf{Train.\ Cost} \\
				\midrule
				LSTM~\cite{hochreiter1997lstm} & Yes & No  & No  & No  & High    \\
				GRU~\cite{cho2014gru}                & Yes & No  & No  & No  & High    \\
				Transformer~\cite{vaswani2017attention}        & No  & Yes & No  & Yes & Medium  \\
				CNN-BiGRU~\cite{cnnbigru2022}          & Yes & No  & No  & No  & High    \\
				Standard ESN~\cite{jaeger2001}       & No  & No  & Yes & No  & Low     \\
				DeepESN~\cite{gallicchio2017}            & No  & No  & Yes & No  & Low     \\
				CNN-ESN~\cite{cesn2024}            & Yes & No  & Yes & No  & Medium  \\
				LSTM-ESN~\cite{lopez2018lstmesn}           & Yes & No  & Yes & No  & Medium  \\
				\textbf{T-ESN (ours)}
				& \textbf{No} & \textbf{Yes} & \textbf{Yes}
				& \textbf{Yes}$^{\ast}$ & \textbf{Low--Med} \\
				\bottomrule
				\multicolumn{6}{l}{\scriptsize
					$^{\ast}$Quadratic in sequence length $T$; mitigated by the short
					sequence design adopted in this work
					(Section~\ref{subsec:training}).}
		\end{tabular}}
	\end{table}
	
	\subsection{Experimental Setup}
	
	\subsubsection{Data Preprocessing}
	\label{subsec:preprocessing}
	
	Raw records were audited for quality issues before training. Features
	with infinite values---\textit{Flow Bytes/s} and \textit{Flow
		Packets/s}, undefined for zero-duration flows---were replaced with NaN
	and discarded, reducing the feature set to 76 dimensions. Duplicate
	records were removed to prevent data leakage. Min-Max normalisation
	was then applied to each remaining feature:
	
	\begin{equation}
		u'_j = \frac{u_j - u_j^{\min}}{u_j^{\max} - u_j^{\min}}
		\label{eq:minmax}
	\end{equation}
	
	This maps all values to $[0,1]$, so features with disparate dynamic
	ranges contribute comparably to both the reservoir dynamics
	(Eq.~\eqref{eq:esn_update}) and the Transformer input projection
	(Eq.~\eqref{eq:proj}). Scaling parameters are computed on the training
	set and applied to the test set to prevent leakage. A stratified
	80:20 train-test split preserved the original class distribution
	across both subsets, consistent with prior
	work~\cite{cnnbigru2022,kurniabudi2020}. Class labels were
	integer-encoded for cross-entropy loss computation.
	
	\subsubsection{Evaluation Metrics}
	
	Four complementary metrics were used given the severe class imbalance.
	Overall accuracy measures the fraction of correctly classified
	samples. Macro-averaged precision, recall, and F1-score weight each
	class equally regardless of sample count:
	
	\begin{equation}
		\mathrm{Macro\text{-}F1} = \frac{1}{C}\sum_{c=1}^{C} F1_c,
		\quad
		F1_c = \frac{2\,P_c\,R_c}{P_c + R_c}
		\label{eq:macrof1}
	\end{equation}
	
	where $C = 14$. Equal per-class weighting makes macro-F1 particularly
	sensitive to minority-class performance---critical for intrusion
	detection, where missing even a single rare attack such as
	Infiltration is a serious operational failure. All models are
	assessed under this identical four-metric protocol.
	
	\subsubsection{Training Configuration}
	\label{subsec:training}
	
	All neural network models were trained under a unified configuration,
	listed in full in Table~\ref{tab:config}.
	
	\begin{table}[!htbp]
		\centering
		\caption{Training configuration for the proposed Transformer-ESN.}
		\label{tab:config}
		\renewcommand{\arraystretch}{0.92}
		\setlength{\tabcolsep}{4pt}
		\begin{minipage}{0.49\linewidth}
			\centering
			\resizebox{\linewidth}{!}{%
				\begin{tabular}{ll}
					\toprule
					\textbf{Parameter} & \textbf{Value} \\
					\midrule
					\multicolumn{2}{l}{\textit{Optimisation}} \\
					Optimiser                       & AdamW \\
					Learning Rate                   & $3\times10^{-4}$ \\
					Weight Decay                    & $1\times10^{-2}$ \\
					Epochs                          & 60 \\
					Batch Size                      & 4096 \\
					Loss Function                   & Weighted Cross-Entropy \\
					Label Smoothing ($\varepsilon$) & 0.05 \\
					Early Stopping Patience         & 15 epochs \\
					Gradient Clipping               & Max norm 1.0 \\
					Class Weights                   & $1/\sqrt{n_c}$, clamped $[0.1,\;10]$ \\
					Train/Test Split                & 80:20 (stratified) \\
					\midrule
					\multicolumn{2}{l}{\textit{ESN Reservoir}} \\
					Reservoir Size ($N$)            & 256 units \\
					Spectral Radius ($r$)           & 0.9 \\
					Input Scaling ($\sigma$)        & 0.1 \\
					Leaking Rate ($\alpha$)         & 1.0 \\
					Connectivity                    & 10\% sparse \\
					\bottomrule
			\end{tabular}}
		\end{minipage}%
		\hfill
		\begin{minipage}{0.49\linewidth}
			\centering
			\resizebox{\linewidth}{!}{%
				\begin{tabular}{ll}
					\toprule
					\textbf{Parameter} & \textbf{Value} \\
					\midrule
					\multicolumn{2}{l}{\textit{Transformer Encoder}} \\
					Attention Heads ($H$)           & 4 \\
					Encoder Layers ($L$)            & 2 \\
					Model Dimension ($d_{\mathrm{model}}$) & 128 \\
					Feed-Forward Dim ($d_{\mathrm{ff}}$)   & 256 \\
					Dropout                         & 0.1 \\
					\midrule
					\multicolumn{2}{l}{\textit{Classification Head}} \\
					Hidden Size                     & 256 \\
					Output Classes ($C$)            & 14 \\
					Total Trainable Parameters      & ${\approx}508\,\mathrm{K}$ \\
					\midrule
					\multicolumn{2}{l}{\textit{Hardware / Software}} \\
					CPU                             & Intel Core i7-12700H \\
					GPU                             & NVIDIA RTX 3060 (6\,GB VRAM) \\
					RAM                             & 32\,GB DDR5 \\
					Python                          & 3.10.12 \\
					PyTorch                         & 2.1.0 \\
					CUDA                            & 11.8 \\
					scikit-learn                    & 1.3.2 \\
					Random Seed                     & 42 \\
					\bottomrule
			\end{tabular}}
		\end{minipage}
	\end{table}
	
	\subsubsection{Random Forest Baseline Configuration}
	\label{subsec:rf_config}
	
	The Random Forest baseline used scikit-learn~1.3.2 under the same
	stratified 80:20 split and preprocessing pipeline
	(Section~\ref{subsec:preprocessing}). Full hyperparameters are given
	in Table~\ref{tab:rf_config} for reproducibility.
	
	\begin{table}[!htbp]
		\centering
		\caption{Random Forest baseline configuration (scikit-learn 1.3.2).}
		\label{tab:rf_config}
		\renewcommand{\arraystretch}{0.9}
		\setlength{\tabcolsep}{5pt}
		\begin{tabularx}{\linewidth}{l X}
			\toprule
			\textbf{Parameter} & \textbf{Value} \\
			\midrule
			n\_estimators       & 100 \\
			max\_depth          & None \\
			min\_samples\_split & 2 \\
			min\_samples\_leaf  & 1 \\
			max\_features       & \texttt{sqrt} \\
			class\_weight       & \texttt{balanced} \\
			bootstrap           & True \\
			random\_state       & 42 \\
			n\_jobs             & $-1$ \\
			\bottomrule
		\end{tabularx}
	\end{table}
	
	All reported metrics represent a single experimental run under the
	described configuration. Future work will include repeated runs with
	varied random seeds and bootstrapped confidence intervals to estimate
	performance variance across both well-supported and ultra-rare
	traffic classes.
	% =========================================================
	% 4. RESULTS
	% =========================================================
	\section{Results}
	\label{sec:results}
	
	\subsection{Overall Performance}
	\label{subsec:overall}
	
	Table~\ref{tab:overall} reports aggregate performance for both models
	evaluated under identical preprocessing and split conditions.
	Transformer-ESN achieves 99.38\% accuracy, 85.34\% macro-precision,
	82.79\% macro-recall, and 73.55\% macro-F1. The reproduced Random Forest
	achieves 99.80\% accuracy and 88.40\% macro-F1.
	
	The accuracy gap is modest (0.42 percentage points). The macro-F1 gap
	is substantial (14.85 percentage points). Section~\ref{subsec:decompose}
	attributes this gap to two separable factors: statistical instability on
	ultra-rare classes, and the structural mismatch between sequential
	inductive biases and pre-aggregated tabular features.
	
	\begin{table}[!htbp]
		\centering
		\caption{Aggregate performance on CIC-IDS2017 (80:20 stratified split).
			Both models reproduced and evaluated under identical conditions.}
		\label{tab:overall}
		\resizebox{\linewidth}{!}{%
			\begin{tabular}{lcccc}
				\toprule
				\textbf{Method}
				& \textbf{Acc.\ (\%)}
				& \textbf{Macro-Prec.\ (\%)}
				& \textbf{Macro-Rec.\ (\%)}
				& \textbf{Macro-F1\ (\%)} \\
				\midrule
				Random Forest (reproduced)~\cite{rf2020}          & 99.80 & 98.20 & 87.30 & 88.40 \\
				\textbf{Transformer-ESN (proposed)} & \textbf{99.38} & \textbf{85.34} & \textbf{82.79} & \textbf{73.55} \\
				\bottomrule
		\end{tabular}}
	\end{table}
	
	\subsection{Comparison with Published Literature}
	\label{subsec:litcomp}
	
	Table~\ref{tab:litcomp} situates the proposed architecture within the
	published landscape of CIC-IDS2017 evaluations. Two caveats govern
	interpretation. First, \emph{direct numerical comparison is not valid}:
	published results vary across preprocessing choices, split ratios, label
	encodings, and class-subset selections~\cite{thakkar2020}. The figures
	are included for contextual orientation, not for head-to-head performance
	ranking. Second, most prior works report aggregate accuracy only; per-class
	and macro-averaged breakdowns are the exception, not the norm.
	
	\begin{table*}[!htbp]
		\centering
		\caption{Comparison of proposed Transformer-ESN against representative
			published results on CIC-IDS2017. ``---'' denotes a metric not reported
			in the original paper. Accuracy and F1 figures are taken directly from
			the cited sources; protocols differ across studies and direct numerical
			comparison is not valid. The proposed model is the only entry with a
			full per-class breakdown reported.}
		\label{tab:litcomp}
		\resizebox{\textwidth}{!}{%
			\begin{tabular}{L{4.0cm} L{3.8cm} C{2.0cm} C{2.0cm} C{2.0cm} C{2.0cm} C{1.6cm}}
				\toprule
				\textbf{Reference}
				& \textbf{Architecture}
				& \textbf{Acc.\ (\%)}
				& \textbf{Prec.\ (\%)}
				& \textbf{Rec.\ (\%)}
				& \textbf{Macro-F1\ (\%)}
				& \textbf{Per-class} \\
				\midrule
				\multicolumn{7}{l}{\textit{Classical Machine Learning}} \\
				Pranto et al.~\cite{kurniabudi2020}
				& Random Forest + SMOTE  & 99.50 & --- & --- & --- & No \\
				Azam et al.~\cite{xgb2023}
				& XGBoost                & 99.60 & --- & --- & --- & No \\
				Ahmad et al.~\cite{svmelm2018}
				& SVM-ELM                & 99.67 & --- & --- & --- & No \\
				\midrule
				\multicolumn{7}{l}{\textit{Deep Recurrent / Convolutional}} \\
				Tang et al.~\cite{grurnn2018}
				& GRU-RNN                & 99.00 & --- & --- & --- & No \\
				Wang et al.~\cite{cnnbigru2022}
				& CNN-BiGRU              & 99.70 & --- & --- & --- & No \\
				Kaur et al.~\cite{cnnlstm2024}
				& CNN-LSTM               & 96.21 & --- & --- & --- & No \\
				\midrule
				\multicolumn{7}{l}{\textit{Reservoir Computing}} \\
				Gallicchio et al.~\cite{gallicchio2017}
				& DeepESN                & 97.60 & --- & --- & --- & No \\
				\midrule
				\multicolumn{7}{l}{\textit{Attention-Based}} \\
				Wu et al.~\cite{rtids2022}
				& Transformer            & 99.17 & --- & --- & --- & No \\
				Zhang et al.~\cite{bilstm_attn2023}
				& BiLSTM + Attention     & 99.40 & --- & --- & --- & No \\
				Li et al.~\cite{trbma2025}
				& TRBMA                  & 98.66 & --- & --- & --- & No \\
				\midrule
				\multicolumn{7}{l}{\textit{Reproduced Baselines (this work)}} \\
				RF (reproduced)
				& Random Forest          & 99.80 & --- & --- & 88.40 & No \\
				\textbf{Proposed}
				& \textbf{Transformer-ESN} & \textbf{99.38} & \textbf{85.34} & \textbf{82.79} & \textbf{73.55} & \textbf{Yes} \\
				\bottomrule
				\multicolumn{7}{l}{%
					\scriptsize $^*$~TRBMA~\cite{trbma2025} reports 98.66\% accuracy and
					prioritises macro-recall over precision by design (BS-OSS hybrid
					sampling); it does not report precision, recall, or macro-F1, so
					those cells are left blank rather than estimated.}
		\end{tabular}}
	\end{table*}
	
	Three observations follow from Table~\ref{tab:litcomp}.
	
	\textbf{Accuracy is within the published range.}
	Transformer-ESN's 99.38\% falls within the 96.21\%--99.70\% range
	reported across all paradigms. It outperforms the standalone Transformer
	of Wu et al.~\cite{rtids2022} (99.17\%) and TRBMA~\cite{trbma2025}
	(98.66\%), and is marginally below CNN-BiGRU (99.70\%, the top of the
	range) and the top classical baselines.
	
	\textbf{Macro-F1 comparison is impossible for most prior work.}
	Ten of the eleven prior entries do not report macro-F1. The proposed
	model is one of the few entries in this literature with a complete
	macro-averaged and per-class breakdown.
	
	\textbf{The evaluation gap is systematic.}
	Every prior entry reports no per-class breakdown. It is therefore
	impossible to determine from the existing literature whether any
	published IDS architecture achieves strong minority-class detection,
	or whether high aggregate scores are driven exclusively by majority-class
	performance.
	
	\subsection{Per-Class Performance}
	\label{subsec:perclass}
	
	Table~\ref{tab:perclass} reports per-class precision, recall, and
	F1-score for the proposed Transformer-ESN across all fourteen traffic
	classes. This is the central result of the paper. Aggregate accuracy
	is insufficient to characterise model behaviour under the severe
	class imbalance of CIC-IDS2017; per-class analysis is required to
	determine where model capacity is actually concentrated.
	
	\begin{table}[!htbp]
		\centering
		\caption{Per-class performance of the proposed Transformer-ESN on the
			CIC-IDS2017 test set. Classes marked $\dagger$ have fewer than 393 test
			samples; per-class metrics for these classes are statistically unreliable
			and should be interpreted as indicative only. The macro average is
			computed over all 14 classes with equal weighting.}
		\label{tab:perclass}
		\resizebox{\linewidth}{!}{%
			\begin{tabular}{L{4.0cm} C{1.4cm} C{1.4cm} C{1.4cm} C{1.4cm}}
				\toprule
				\textbf{Class}
				& \textbf{Prec.}
				& \textbf{Rec.}
				& \textbf{F1}
				& \textbf{Support} \\
				\midrule
				\multicolumn{5}{l}{\textit{Well-supported classes ($\geq$393 samples)}} \\
				Benign                  & 1.00 & 0.99 & 1.00 & 454{,}620 \\
				DDoS                    & 1.00 & 1.00 & 1.00 &  25{,}606 \\
				DoS GoldenEye           & 0.97 & 0.99 & 0.98 &   2{,}059 \\
				DoS Hulk                & 0.97 & 1.00 & 0.98 &  46{,}215 \\
				DoS Slowhttptest        & 0.90 & 0.99 & 0.94 &   1{,}100 \\
				DoS Slowloris           & 0.99 & 0.99 & 0.99 &   1{,}159 \\
				FTP-Patator             & 0.99 & 0.99 & 0.99 &   1{,}588 \\
				PortScan                & 0.99 & 1.00 & 1.00 &  31{,}786 \\
				SSH-Patator             & 0.96 & 0.98 & 0.97 &   1{,}179 \\
				Botnet                  & 0.58 & 0.65 & 0.61 &     393   \\
				\midrule
				\multicolumn{5}{l}{\textit{Ultra-rare classes (fewer than 393 test samples)}} \\
				WA -- Brute Force$\dagger$ & 0.20 & 0.94 & 0.33 & 301 \\
				WA -- XSS$\dagger$          & 1.00 & 0.03 & 0.06 & 130 \\
				WA -- SQL Injection$\dagger$& 0.09 & 0.25 & 0.13 &   4 \\
				Infiltration$\dagger$       & 0.33 & 0.29 & 0.31 &   7 \\
				\midrule
				\textbf{Macro Average}  & \textbf{0.85} & \textbf{0.83} & \textbf{0.74} & 566{,}147 \\
				\bottomrule
		\end{tabular}}
	\end{table}
	
	Figure~\ref{fig:f1bar} visualises the F1-score distribution across
	all fourteen classes, and Figure~\ref{fig:four_metric} in
	Section~\ref{subsec:decompose} presents a four-metric comparison
	between Transformer-ESN and Random Forest.
	
	\subsection{Structural Analysis of Results}
	\label{subsec:structural}
	
	Three structural observations emerge from Table~\ref{tab:perclass}
	and Figure~\ref{fig:f1bar}.
	
	\textbf{Well-supported classes perform strongly.}
	All ten classes with test support of 393 samples or more achieve
	F1-scores of 0.61 or above. Nine of these ten exceed 0.94, including
	all four DoS variants (Hulk, GoldenEye, Slowloris, Slowhttptest),
	both credential-attack classes (FTP-Patator, SSH-Patator), PortScan,
	DDoS, and Benign. The model demonstrates reliable discrimination
	wherever training signal is sufficient.
	
	\textbf{Ultra-rare classes dominate macro-F1 degradation.}
	The four classes marked $\dagger$ (Web Attack--Brute Force, XSS, SQL
	Injection, Infiltration) have test-set support of 301, 130, 4, and 7
	samples respectively. Their F1-scores (0.06--0.33) drive the macro
	average from approximately 97\% (computed over the ten well-supported
	classes alone) down to 73.55\% across all fourteen. A single
	misclassification in the 4-sample SQL Injection test set constitutes
	a 25\% error rate for that class. Per-class metrics at these support
	levels are not statistically stable estimators of architectural
	capability; individual prediction outcomes dominate the reported
	values, regardless of architecture.
	
	\textbf{Botnet is a genuine intermediate case.}
	With 393 test samples and F1 of 0.61, Botnet occupies an intermediate
	position. Its precision (0.58) and recall (0.65) indicate that the
	model achieves meaningful but incomplete separation---a pattern
	consistent with the known feature overlap between Botnet ARES traffic
	and background benign flows in CIC-IDS2017~\cite{cicids2017}. This
	represents a genuine detection challenge distinct from the statistical
	instability affecting the ultra-rare classes, and constitutes a
	concrete target for future improvement.
	
	\subsection{Web Attack -- XSS: Precision-Recall Inversion}
	\label{subsec:xss}
	
	Web Attack -- XSS exhibits a qualitatively distinct failure mode:
	precision 1.00 but recall 0.03, yielding F1 of 0.06. This
	precision-recall inversion indicates that the model has learned
	to assign the XSS label extremely conservatively, predicting it
	correctly for only the most unambiguous examples while
	misclassifying the majority as a more common class.
	
	With 130 test samples, approximately 126 XSS flows are being
	misclassified. The inverse pattern in Web Attack -- Brute Force
	(recall 0.94, precision 0.20) strongly suggests that Brute Force
	absorbs the majority of misclassified XSS instances: the two
	classes share training feature distributions. This is consistent
	with their common origin---both are HTTP-based application-layer
	attacks with overlapping flow-level statistics.
	
	This failure mode is not a statistical artefact of small support; it
	is a real classification failure attributable to insufficient training
	signal (522 XSS training samples) combined with feature overlap with
	a semantically adjacent but more numerous class. No architectural
	modification within this feature representation can fully resolve it
	without additional XSS-specific training data.
	
	% =========================================================
	% DISCUSSION
	% =========================================================
	\section{Discussion}
	\label{sec:discussion}
	
	\subsection{Decomposing the Macro-F1 Gap}
	\label{subsec:decompose}
	
	The 14.85 percentage point macro-F1 gap between Transformer-ESN
	(73.55\%) and Random Forest (88.40\%) is the central result requiring
	explanation. Our analysis suggests that the observed macro-F1 gap is
	largely explained by two major factors. Treating these as a single
	undifferentiated cause risks an inaccurate characterisation of the
	architecture's capability, and each is therefore considered separately
	below.
	
	\textbf{Factor 1: Statistical unreliability on ultra-rare classes.}
	Four test classes have support of 4, 7, 130, and 301 samples. At these
	support levels, per-class metrics are unlikely to be statistically
	stable for any classifier: a single correct or incorrect prediction
	changes SQL Injection's F1 by 25\% and Infiltration's F1 by 14\%. Under
	these conditions, the macro average is a less reliable indicator of
	architectural capability, since it is highly sensitive to individual
	prediction outcomes. This instability appears to be a property of the
	dataset and its evaluation protocol rather than of the architecture
	under test.
	
	\textbf{Factor 2: Structural mismatch between architecture and data.}
	CICFlowMeter produces pre-aggregated statistical summaries over completed
	flows~\cite{cicids2017}; each row represents an independent observation
	and does not encode an ordered position within a packet-level sequence.
	Sequential inductive biases, such as those imposed by an ESN reservoir
	and a Transformer encoder, are less likely to provide substantial
	benefit on pre-aggregated tabular flow features and may introduce
	unnecessary modelling complexity relative to the information actually
	present in the data. Random Forest, in contrast, imposes no such
	assumption: it treats each row as an independent point in feature
	space, an approach that appears well aligned with the tabular nature
	of the input. This structural alignment is likely a major contributing
	factor to the observed macro-F1 gap, rather than evidence of a
	difference in general representational capacity between the two
	approaches.
	
	Table~\ref{tab:decompose} presents an approximate decomposition of
	these two factors. The subset values were obtained by recomputing
	macro-F1 over progressively restricted class subsets and are therefore
	analytical estimates rather than independently measured quantities;
	they should be interpreted as indicative of relative magnitude rather
	than as precise, additive contributions.
	
	\begin{table}[!htbp]
		\centering
		\caption{Macro-F1 computed over class subsets to isolate the
			approximate contribution of ultra-rare classes to the observed
			performance gap. Values are analytical estimates derived by recomputing
			macro-F1 over successive class subsets, not independently measured
			quantities. The structural mismatch factor remains relevant for Botnet
			even within the well-supported group.}
		\label{tab:decompose}
		\renewcommand{\arraystretch}{1.15}
		\begin{tabular}{L{4.6cm} C{2.0cm}}
			\toprule
			\textbf{Class Subset} & \textbf{Macro-F1 (\%)} \\
			\midrule
			All 14 classes (reported)    & 73.55 \\
			10 well-supported classes only & $\approx$96.40 \\
			9 high-performing classes (excl.\ Botnet) & $\approx$98.30 \\
			\midrule
			\textit{Gap attributable to ultra-rare classes} & $\approx$22.85 pp \\
			\textit{Gap attributable to Botnet}              & $\approx$1.90 pp \\
			\bottomrule
			\multicolumn{2}{l}{\scriptsize pp = percentage points; approximate, see text}
		\end{tabular}
	\end{table}
	
	\subsection{Botnet as the Boundary Case}
	\label{subsec:botnet}
	
	With 393 test samples, the Botnet class has a sample size sufficient to
	produce statistically meaningful per-class estimates, distinguishing
	it from the ultra-rare classes discussed above. The partial separation
	achieved---precision 0.58, recall 0.65---suggests that the architecture
	captures a genuine learning signal for this class, although separation
	remains incomplete. Botnet ARES traffic in CIC-IDS2017 is known to
	exhibit packet-timing and flow-length statistics that partially overlap
	with normal background traffic~\cite{cicids2017}, which is likely to
	make it intrinsically harder to separate than volumetric attack
	categories. The Botnet result indicates that the architecture's
	limitations are not uniform across all minority classes: moderate-support
	classes such as Botnet appear to remain learnable even when ultra-rare
	classes dominate the overall macro-F1 penalty. This distinction
	positions Botnet as a tractable improvement target, separate from the
	structurally unreliable ultra-rare class estimates. Without a published
	per-class breakdown for the reproduced Random Forest baseline, whether
	Random Forest achieves comparable separation on this specific class
	remains an open question, and the comparison at the Botnet level should
	be treated as inconclusive rather than as a demonstrated advantage for
	either method.
	
	\subsection{Where the Architecture Provides Value}
	\label{subsec:value}
	
	While the aggregate macro-F1 comparison favours Random Forest on this
	dataset, several architectural properties of Transformer-ESN are not
	captured by that single metric and merit separate consideration.
	
	\textbf{Training stability.}
	The fixed ESN reservoir removes the need for recurrent backpropagation,
	confining gradient computation to the Transformer encoder and
	classification head. This design is expected to substantially reduce
	the effects of vanishing and exploding gradients that constrain fully
	trained LSTM and GRU architectures, likely reducing sensitivity to
	initialisation and learning-rate choices relative to end-to-end
	recurrent training~\cite{grurnn2018}. This is broadly consistent with
	recent findings that ESN-based reservoirs mitigate the gradient
	instability associated with backpropagation-through-time in LSTM/GRU
	architectures~\cite{alam2026drlesn}, though that result was obtained on
	NetFlow-based IoT datasets rather than CIC-IDS2017 and should not be
	read as CIC-IDS2017-specific evidence.
	
	\textbf{Parameter efficiency.}
	The 508K-parameter trainable footprint is comparatively small relative
	to attention-based IDS architectures more broadly, since the model
	avoids full end-to-end training of its recurrent component. As
	summarised in Table~\ref{tab:architecture_compare}, Transformer-ESN is
	the only reviewed architecture combining a gradient-free reservoir with
	an attention mechanism, which structurally limits the fraction of total
	model capacity that requires gradient updates. The
	$256{\times}76$ input matrix and $256{\times}256$ recurrent weight
	matrix are fixed at initialisation and do not contribute to gradient
	computation during training. A precise quantitative comparison against
	specific published attention-based IDS models would require parameter
	counts not consistently reported across that literature, and is
	therefore left as a direction for future benchmarking rather than
	claimed here.
	
	\textbf{Validated discrimination on well-supported classes.}
	On the nine classes achieving F1 above 0.94, the reservoir-to-attention
	pipeline demonstrates discriminative behaviour consistent with
	sufficient training data being available. This result supports the
	architectural design in settings where the sequential inductive bias
	is appropriate, such as packet-level stream data.
	
	\textbf{Structural readiness for sequential data.}
	The architecture is designed for settings in which the ESN reservoir
	can exploit genuine sequential structure in the input. On per-flow
	tabular features, this structure is largely absent, meaning the
	architecture is likely operating outside its intended application
	domain in the present evaluation. Taken together with the findings
	above, this suggests a mismatch between the architecture's intended
	operating regime and the feature representation available in
	CIC-IDS2017, rather than a fundamental limitation of the architecture
	itself.
	
	These architectural considerations point toward specific deployment
	contexts in which the properties described above are expected to
	translate into practical advantages, as discussed next.
	
	\subsection{Practical Deployment Scenarios}
	\label{subsec:deployment}
	
	The structural properties of Transformer-ESN have concrete implications
	for deployment contexts that aggregate accuracy comparisons do not
	fully capture. Three scenarios illustrate where the architecture is
	expected to provide practical value relative to classical tabular
	methods.
	
	\textbf{Packet-level and session-stream monitoring.}
	Enterprise NIDS deployments increasingly operate on raw packet captures
	or reassembled TCP session streams rather than pre-aggregated flow
	statistics~\cite{buczak2016survey,khraisat2019survey}. In these
	settings, each input is a genuine temporal sequence: packet inter-arrival
	times, payload length progressions, and flag sequences carry information
	that is discarded by per-flow aggregation. The ESN reservoir is
	designed to exploit this type of structure: its recurrent dynamics
	maintain a fading memory of recent inputs, and its spectral radius
	governs the effective temporal horizon over which past context
	influences the current state. When operating on such data, the fixed
	reservoir is expected to provide a computationally inexpensive nonlinear
	summary of recent traffic history, while the Transformer encoder
	attends to dependencies across multiple time steps within a session.
	Random Forest, which treats each observation as independent, is
	unlikely to model this cross-packet temporal structure regardless of
	feature engineering effort.
	
	\textbf{Resource-constrained and edge deployment.}
	Network monitoring at the edge---on switches, routers, or embedded
	security appliances---imposes strict constraints on model size,
	retraining frequency, and labelled data availability. The
	Transformer-ESN's 508K trainable parameter footprint is comparatively
	small relative to attention-based IDS architectures more broadly, and
	its fixed reservoir substantially reduces the computational burden
	associated with recurrent training. In environments where labelled
	attack traffic is scarce and retraining must be completed within
	strict time windows, this training efficiency and compact parameter
	count represent practical advantages that accuracy figures alone do
	not quantify.
	
	\textbf{Federated and collaborative monitoring.}
	Multi-organisational threat intelligence sharing is constrained by
	data privacy regulations that prohibit raw traffic exchange. In a
	federated NIDS deployment, only the trainable encoder weights (508K
	parameters) would need to be aggregated across nodes, while each
	participant retains its fixed reservoir locally. This is expected to
	create substantially lower communication overhead than federating
	fully trained recurrent architectures, and the reservoir's fixed random
	initialisation allows participants to share an identical reservoir
	without coordination. These properties suggest Transformer-ESN as a
	plausible candidate for privacy-preserving collaborative
	detection~\cite{idrissi2023fedanids,feddef2023}.
	
	\subsection{Architecture Selection Guidelines}
	\label{subsec:selection}
	
	The experimental results, taken together with the deployment
	considerations above, suggest a practical decision boundary between
	the proposed architecture and classical tabular methods for
	practitioners selecting a classifier for a given deployment context.
	
	Random Forest appears to be the more appropriate choice when the input
	consists of pre-aggregated per-flow statistics (as in CIC-IDS2017),
	labelled data is abundant, and the operational requirement is
	maximising per-class accuracy across all categories, including
	ultra-rare ones. Its structural alignment with row-wise tabular
	features, combined with the native ensemble averaging that reduces
	variance on minority classes when sufficient training samples exist,
	makes it difficult to surpass on this specific input representation.
	The macro-F1 advantage observed here (88.40\% vs.\ 73.55\%) appears to
	be a direct consequence of this alignment and should not be interpreted
	as evidence of superior general representational capacity.
	
	Transformer-ESN is likely to become the preferable choice under three
	conditions: (i) the input consists of ordered sequences rather than
	independent per-flow summaries---packet captures, NetFlow sequences,
	or HTTP session logs where the temporal ordering of records encodes
	threat behaviour not captured by any single-row statistic; (ii)
	training stability is a hard constraint, as the fixed reservoir is
	expected to substantially reduce the hyperparameter sensitivity and
	gradient instability that affect LSTM and GRU training on imbalanced
	security data; or (iii) the deployment requires incremental or
	federated model updates, where the compact trainable component offers
	communication and computation efficiency that tree ensembles are
	unlikely to match. The architecture is not proposed as a universal
	replacement for Random Forest on tabular IDS benchmarks, but rather as
	a structurally motivated choice for the sequential and distributed
	monitoring scenarios that increasingly characterise modern network
	security infrastructure.
	
	Table~\ref{tab:architecture_compare} summarises these properties
	relative to the Random Forest baseline and representative prior work.
	
	\begin{table}[!htbp]
		\centering
		\caption{Qualitative comparison of architectural properties. Checkmarks
			denote property presence; the proposed Transformer-ESN is the only entry
			with all four properties simultaneously.}
		\label{tab:architecture_compare}
		\renewcommand{\arraystretch}{1.15}
		\resizebox{\linewidth}{!}{%
			\begin{tabular}{L{3.6cm} C{1.6cm} C{1.6cm} C{1.6cm} C{1.8cm}}
				\toprule
				\textbf{Architecture}
				& \textbf{Temporal encoding}
				& \textbf{Gradient-free reservoir}
				& \textbf{Attention mechanism}
				& \textbf{Per-class eval.} \\
				\midrule
				Random Forest          & \xmark & \xmark & \xmark & \xmark \\
				GRU-RNN~\cite{grurnn2018}
				& \cmark & \xmark & \xmark & \xmark \\
				CNN-BiGRU~\cite{cnnbigru2022}
				& \cmark & \xmark & \xmark & \xmark \\
				DeepESN~\cite{gallicchio2017}
				& \cmark & \cmark & \xmark & \xmark \\
				Transformer~\cite{rtids2022}
				& \xmark & \xmark & \cmark & \xmark \\
				TRBMA~\cite{trbma2025}
				& \cmark & \xmark & \cmark & \xmark \\
				\textbf{Transformer-ESN (proposed)}
				& \cmark & \cmark & \cmark & \cmark \\
				\bottomrule
		\end{tabular}}
	\end{table}
	
	Figure~\ref{fig:crossparadigm} compares average reported accuracy
	across these paradigms, while Figure~\ref{fig:evolution_ids} traces
	the paradigm progression reviewed in Section~\ref{sec:litreview}
	that motivates the proposed coupling.
	
	\begin{figure}[!htbp]
		\centering
		\safefig{Figure5_Average_Architecture_Accuracy.png}{Average accuracy per learning paradigm}
		\caption{Average reported accuracy on CIC-IDS2017 for the best representative model within each of the four learning paradigms reviewed. Transformer-ESN (99.38\%) is competitive with the strongest deep recurrent/convolutional baseline (CNN-BiGRU~\cite{cnnbigru2022}, 99.70\%) and the strongest classical baseline (SVM-ELM~\cite{svmelm2018}, 99.67\%), and outperforms the strongest standalone attention-based model (Transformer~\cite{rtids2022}, 99.17\%). As with Figure~\ref{fig:accuracy_comparison}, cross-study comparisons are contextual rather than statistically controlled, and the proposed model remains the only entry with a complete per-class breakdown (Table~\ref{tab:perclass}).}
		\label{fig:crossparadigm}
	\end{figure}
	
	\begin{figure}[!htbp]
		\centering
		\safefig{Figure6_Evolution_IDS.png}{Evolution of intrusion detection paradigms}
		\caption{Timeline summarising the progression of network intrusion detection paradigms reviewed in Section~\ref{sec:litreview}: classical machine learning, deep recurrent/convolutional networks, reservoir computing, and attention-based architectures, culminating in the reservoir-to-attention coupling proposed in this work. The figure is intended as a conceptual overview of paradigm succession rather than a chronologically exhaustive record.}
		\label{fig:evolution_ids}
	\end{figure}
	
	Together, Figure~\ref{fig:crossparadigm} and Figure~\ref{fig:evolution_ids}
	summarise, respectively, where the proposed model sits relative to each
	paradigm's best result and how the four paradigms reviewed in
	Section~\ref{sec:litreview} led to the reservoir-to-attention coupling
	examined in this paper.
	
	\begin{figure}[!htbp]
		\centering
		\safefig{Figure7_Transformer_ESN_Workflow.png}{Transformer-ESN training workflow}
		\caption{Training workflow of the proposed Transformer-ESN, from data loading through convergence checking and final evaluation. Reservoir weights $W_{\mathrm{in}}$ and $W$ are initialised once and frozen throughout; only the Transformer encoder and classification head are updated via backpropagation at each iteration.}
		\label{fig:training_flow}
	\end{figure}
	
	\begin{figure}[!htbp]
		\centering
		\safefig{Figure8_Performance_Summary.png}{Per-class F1 performance summary}
		\caption{Per-class F1-score of the proposed Transformer-ESN on the CIC-IDS2017 test set, summarising the results of Table~\ref{tab:perclass}. Classes marked $\dagger$ have fewer than 393 test samples, so their F1 estimates are statistically unstable. The bimodal pattern---F1~$\geq$~0.94 on nine of ten well-supported classes against near-zero F1 on four ultra-rare classes---accounts for the divergence between overall accuracy (99.38\%) and macro-F1 (73.55\%).}
		\label{fig:f1bar}
	\end{figure}
	
	\subsection{Contextualisation Against Literature}
	\label{subsec:context}
	
	Transformer-ESN's 99.38\% accuracy sits within the 96.21\%--99.70\%
	range reported across all four paradigms reviewed in
	Table~\ref{tab:litcomp}. It outperforms the standalone Transformer
	of Wu et al.~\cite{rtids2022} and TRBMA~\cite{trbma2025} on accuracy,
	and is marginally below CNN-BiGRU~\cite{cnnbigru2022} and the top
	classical results.
	
	Direct numerical comparison with published figures should be treated
	with caution given protocol heterogeneity across studies~\cite{thakkar2020}.
	What can be reasonably stated is that the proposed architecture's
	accuracy does not appear degraded relative to the published range, and
	that its per-class performance on well-represented classes is
	competitive with the strongest published results. The reported
	limitation is specifically on ultra-rare classes---a limitation that is
	plausibly shared by many published architectures, though rarely
	reported, since most published evaluations omit per-class analysis.
	
	Taken together, these results point to a broader conclusion about model
	selection in network intrusion detection. The macro-F1 gap observed in
	this study appears to be largely explained by two identifiable and
	partially separable factors: the statistical instability of ultra-rare
	class estimates, and a structural mismatch between sequential inductive
	biases and the pre-aggregated, non-sequential nature of CICFlowMeter
	features. Random Forest's advantage on this benchmark is best
	understood as a consequence of its structural fit to tabular, row-independent
	data, rather than a general superiority in representational capacity.
	Conversely, Transformer-ESN's design remains well motivated for data
	representations where genuine sequential structure exists, such as
	packet-level captures or session-ordered logs, and its training
	stability and parameter efficiency offer practical benefits in
	resource-constrained or federated settings regardless of the tabular
	benchmark result. The broader implication is that architecture
	selection in intrusion detection should be guided primarily by whether
	the available data representation carries genuine sequential
	information, rather than by benchmark accuracy alone; this framing
	motivates the directions for future work outlined below.
	
	\subsection{Implications for Future Work}
	\label{subsec:implications}
	
	Three actionable directions follow from these findings.
	
	\textbf{Packet-level sequential data.}
	The sequential inductive bias of the Transformer-ESN architecture is
	expected to be most beneficial when applied to data with genuine
	sequential structure. Evaluation on packet-level sequences within flows,
	or on temporally-ordered multi-flow session logs, would provide the
	architectural preconditions for the ESN reservoir and Transformer
	encoder to contribute as intended. This expectation is consistent with
	recent packet-level frameworks showing that flow-independent sequential
	processing improves detection suitability relative to purely flow-level
	aggregation~\cite{hore2024sequential}.
	
	\textbf{Few-shot strategies for ultra-rare classes.}
	The extreme rarity of several CIC-IDS2017 classes (4--36 total samples)
	makes reliable evaluation difficult under standard train-test splits.
	Future work should combine the proposed architecture with explicit
	few-shot learning, data augmentation (e.g., flow-level WGAN-GP
	synthesis~\cite{thakkar2020}), or open-set detection strategies for
	ultra-rare classes, evaluated under protocols designed for low-support
	regimes.
	
	\textbf{Federated and incremental deployment.}
	The fixed gradient-free reservoir creates a natural basis for federated
	learning: only the compact trainable component (508K parameters) would
	need to be updated or aggregated across distributed nodes, while
	reservoir weights remain locally fixed. Fully trainable recurrent
	architectures are unlikely to offer this structural advantage.
	Investigating this property for distributed or edge-deployed IDS is a
	direct extension of this work. This aligns with the broader federated-IDS
	literature, which identifies communication efficiency and reduced
	trainable footprint as central motivations for federated
	adoption~\cite{hernandezramos2025federated}. A recent federated
	deployment tested directly on CIC-IDS2017 among other benchmarks
	reported its lowest accuracy specifically on CIC-IDS2017 (98.37\%,
	versus 99.65\% on a vehicular network dataset), underscoring that
	federated compactness does not by itself guarantee strong performance
	on this benchmark~\cite{almansour2025adaptive}.
	
	% =========================================================
	% 6. CONCLUSION
	% =========================================================
	\section{Conclusion}
	\label{sec:conclusion}
	
	Network intrusion detection remains one of the most demanding
	classification problems in applied machine learning: class
	distributions are severely skewed, minority-class attacks carry
	disproportionate operational risk, and deployment environments
	impose hard constraints on training latency and labelled data
	availability. Existing deep learning approaches address these challenges either
	through purely recurrent architectures---which couple temporal feature
	extraction and discriminative classification into a single
	gradient-dependent pipeline subject to vanishing gradient instability
	and high sensitivity to initialisation---or through Transformer-only
	models applied directly to static feature vectors, which lack an
	explicit mechanism for capturing the nonlinear temporal dynamics
	present in ordered network traffic.
	
	This paper addressed both limitations through a hybrid
	Transformer--Echo State Network (Transformer-ESN)
	architecture in which the two concerns are deliberately separated.
	A fixed, randomly initialised ESN reservoir first maps the
	76-dimensional preprocessed network flow feature space into a
	256-dimensional nonlinear state representation without any
	gradient-based optimisation~\cite{jaeger2001,lukosevicius2012}.
	The resulting reservoir states are then processed by a multi-head
	self-attention Transformer encoder~\cite{vaswani2017attention},
	which learns long-range inter-feature dependencies across the
	projected space before a softmax classification head produces
	per-class probability scores. While prior work has applied Echo
	State Networks to intrusion detection in standalone
	configurations~\cite{esn_dos2021,mlesn2020,cesn2024} and
	Transformer encoders have been widely adopted for network traffic
	classification~\cite{rtids2022,flowformer2024,long2024transformer},
	the specific coupling proposed here---using a fixed ESN reservoir
	as a zero-cost nonlinear projection layer that feeds directly into
	a Transformer encoder for multi-class network intrusion
	detection---has not, to the best of our knowledge, been previously
	reported in the literature.
	
	Evaluated on the CIC-IDS2017 benchmark~\cite{cicids2017},
	the Transformer-ESN architecture achieved 99.38\% overall
	accuracy, competitive with the range of published results on this
	dataset~\cite{kurniabudi2020,xgb2023,svmelm2018,grurnn2018,
		cnnlstm2024,cnnbigru2022,rtids2022,trbma2025}. The total
	trainable footprint is approximately 508K parameters, with the
	remainder of the model consisting entirely of fixed reservoir
	weights that require no gradient computation at any point during
	training or inference. Against a Random Forest baseline~\cite{kurniabudi2020}
	evaluated under identical experimental conditions, the macro-F1 gap
	(73.55\% vs.\ 88.40\%) is explained by two specific and tractable
	factors: statistical unreliability of per-class metrics at
	ultra-low test support (4--301 samples for four classes), and
	architectural mismatch between sequential inductive biases and
	row-wise tabular features that lack genuine sequential order.
	Across the ten well-represented classes, the architecture achieves
	F1-scores of 0.61--1.00, with nine of ten exceeding 0.94, confirming
	that the performance gap is localised entirely to the sparse-support
	regime rather than reflecting a generalised architectural shortcoming.
	
	Beyond raw accuracy, the architecture's most practically relevant
	property is the clean separation between its two stages. Because
	reservoir weights are fixed at initialisation~\cite{jaeger2001,
		scardapane2017}, the model eliminates recurrent backpropagation
	entirely, reducing training complexity and improving convergence
	stability relative to end-to-end recurrent
	baselines~\cite{grurnn2018,yin2017deep}. Only the Transformer
	encoder and classification head carry trainable parameters,
	making the total trainable parameter count compact and
	well-suited to settings where labelled traffic samples are
	scarce and model retraining must be completed under strict time
	constraints. This separation also creates a principled basis for
	incremental and federated learning, since only the compact
	trainable component needs to be updated or aggregated across
	distributed nodes---an important practical advantage that
	purely end-to-end architectures do not share~\cite{buczak2016survey,
		khraisat2019survey}.
	
	The principal contribution of this work is not a performance
	claim but a structural one: the reservoir-to-attention coupling
	is a training-stable, parameter-efficient approach that
	eliminates recurrent backpropagation while preserving the
	capacity for global contextual
	reasoning~\cite{gallicchio2017,vaswani2017attention}. Its
	limitations on CIC-IDS2017 are specific and tractable---
	attributable to feature representation mismatch and dataset
	sparsity~\cite{ring2019survey,thakkar2020} rather than to
	generalised architectural shortcomings. The architecture is most
	appropriately evaluated in settings where genuine sequential
	structure exists in the input, such as packet-level streams or
	temporally-ordered session logs, where its sequential inductive
	bias becomes an advantage rather than a liability.
	
	Taken together, the experimental results and architectural
	analysis establish Transformer-ESN as a competitive, principled,
	and computationally efficient framework for multi-class network
	intrusion detection. The proposed reservoir-to-attention coupling
	represents a complementary direction to existing hybrid
	architectures~\cite{cnnlstm2024,cnnbigru2022,bilstm_attn2023}
	and opens a clear path toward deployable, updatable, and
	interpretable detection systems in operational cybersecurity
	environments.
	
	\subsection{Future Scope}
	
	\subsubsection{Packet-Level and Sequential Data Evaluation}
	
	The structural mismatch identified in this work---between the
	architecture's sequential inductive bias and the row-wise
	tabular nature of CIC-IDS2017 flow features---is expected to
	diminish when the model operates on genuinely ordered inputs.
	Future work should evaluate Transformer-ESN on packet-level
	streams and temporally-ordered session logs, where ESN reservoir
	dynamics and Transformer attention can exploit real temporal
	ordering rather than imposing it artificially on independent
	feature vectors~\cite{grurnn2018,yin2017deep}.
	
	This expectation is supported, with an important qualification,
	by recent work reformulating CIC-IDS2017 as a temporal task under
	controlled padding and split conditions~\cite{moczkodan2026transformers}.
	On genuinely sequential, non-padded input windows, a standalone
	Transformer achieved the highest macro-F1 of any evaluated
	architecture, consistent with the expectation that sequential
	inductive biases become an advantage once real temporal structure
	is present. However, the same study found that padding convention
	and split protocol---not architecture---were the dominant factors
	governing reported performance: naive zero-padding degraded the
	Transformer's macro-F1 by 0.24, and false-alarm rate under a
	leakage-free split increased by a factor of 67 relative to a
	conventional random split. This indicates that any future
	packet-level or session-level evaluation of Transformer-ESN must
	explicitly control for padding scheme and enforce leakage-free
	splitting; absent these controls, reported gains from sequential
	modeling are not distinguishable from evaluation-protocol
	artifacts~\cite{moczkodan2026transformers}.
	
	\subsubsection{Rare-Class Robustness and Evaluation Protocols}
	
	Four attack categories in CIC-IDS2017 had test support of
	4--301 samples, making per-class F1-scores in those categories
	statistically unreliable. Future work should combine few-shot
	learning strategies, targeted augmentation for ultra-rare attack
	classes, and evaluation protocols specifically designed for the
	low-support regime---such as stratified cross-validation with
	bootstrapped confidence intervals---to produce assessments that
	distinguish genuine architectural limitations from artefacts of
	dataset sparsity~\cite{ring2019survey,thakkar2020,zhou2020ensemble}.
	
	\subsubsection{Statistical Rigour and Repeated Experimental Evaluation}
	
	All results in this paper represent a single experimental run under a
	fixed random seed. While this is consistent with a substantial portion
	of published IDS evaluations, it precludes reporting confidence
	intervals or testing for statistical significance of performance
	differences between architectures. Future work should conduct repeated
	runs across a range of random seeds (recommended minimum: 10--30 runs)
	and report mean performance with standard deviation or bootstrapped
	95\% confidence intervals for each metric. For macro-F1 in particular,
	where minority-class outcomes can shift aggregate values substantially
	on single runs, statistical characterisation of variance is a necessary
	condition for drawing reliable conclusions about architectural
	comparisons. Significance testing (e.g., paired Wilcoxon signed-rank
	tests across cross-validation folds) should be applied when comparing
	Transformer-ESN against the Random Forest baseline to determine whether
	the observed macro-F1 gap is statistically reliable or within
	run-to-run variance bounds.
	
	\subsubsection{Federated and Privacy-Preserving IDS}
	
	Collaborative intrusion detection across organisational
	boundaries is constrained by data privacy regulations that
	prohibit raw traffic sharing. Future research should
	investigate federated Transformer-ESN frameworks in
	which multiple stakeholders jointly optimise the trainable
	encoder using privacy-preserving gradient aggregation. The
	compact parameterisation of the trainable component---limited
	to the Transformer encoder and classification head, with
	reservoir weights fixed and never requiring
	transmission---is particularly advantageous for
	communication-efficient federated
	aggregation~\cite{buczak2016survey,khraisat2019survey}.
	
	\subsubsection{Adaptive and Structured Reservoir Design}
	
	The random sparse reservoir topology adopted in this work
	may not optimally represent all classes of network traffic
	dynamics. Future work should systematically evaluate
	structured reservoir topologies---including ring networks,
	small-world graphs, and spectral graph-designed
	reservoirs---alongside adaptive initialisation strategies
	that tune spectral radius and connectivity density to the
	statistical properties of the target traffic
	distribution~\cite{gallicchio2020,gallicchio2017,scardapane2017}.
	Adaptive reservoir designs may prove particularly effective
	for traffic distributions with non-stationary temporal
	structure, where fixed random projections are a limiting
	assumption.
	
	\subsubsection{Explainable Intrusion Detection}
	
	Operational security analysts require human-interpretable
	justifications for automated classification decisions before
	acting on model outputs. Future work should integrate
	attention heatmap visualisation and feature attribution
	methods to identify which reservoir states---and by
	extension, which temporal traffic patterns---are most
	predictive of specific attack categories. The multi-head
	attention weight matrices produced by the Transformer
	encoder offer a natural and already-available
	interpretability mechanism that can be leveraged without
	adding architectural
	overhead~\cite{rtids2022,bilstm_attn2023,liu2020hierarchical}.
	
	
	% =========================================================
	% REFERENCES
	% =========================================================
	% =========================================================
	% REFERENCES (fixed: citation-order matched, 10 uncited entries removed)
	% =========================================================
	% =========================================================
	% REFERENCES (reordered by first citation appearance)
	% =========================================================
	\begin{thebibliography}{99}
		
		\bibitem{buczak2016survey}
		A.L.\ Buczak and E.\ Guven,
		``A Survey of Data Mining and Machine Learning Methods for Cyber Security Intrusion Detection,''
		IEEE Communications Surveys \& Tutorials,
		vol.\ 18,
		no.\ 2,
		pp.\ 1153--1176,
		2016.
		
		\bibitem{khraisat2019survey}
		A.\ Khraisat and I.\ Gondal and P.\ Vamplew and J.\ Kamruzzaman,
		``Survey of Intrusion Detection Systems: Techniques, Datasets and Challenges,''
		Cybersecurity,
		vol.\ 2,
		no.\ 1,
		pp.\ 1--22,
		2019.
		
		\bibitem{denning1987}
		D.E.\ Denning,
		``An Intrusion-Detection Model,''
		IEEE Transactions on Software Engineering,
		vol.\ SE-13,
		no.\ 2,
		pp.\ 222--232,
		1987.
		
		\bibitem{vinayakumar2019deep}
		R.\ Vinayakumar and M.\ Alazab and K.P.\ Soman and P.\ Poornachandran and A.\ Al-Nemrat and S.\ Venkatraman,
		``Deep Learning Approach for Intelligent Intrusion Detection System,''
		IEEE Access,
		vol.\ 7,
		pp.\ 41525--41550,
		2019.
		
		\bibitem{yin2017deep}
		C.\ Yin and Y.\ Zhu and J.\ Fei and X.\ He,
		``A Deep Learning Approach for Intrusion Detection Using Recurrent Neural Networks,''
		IEEE Access,
		vol.\ 5,
		pp.\ 21954--21961,
		2017.
		
		\bibitem{hochreiter1997lstm}
		S.\ Hochreiter and J.\ Schmidhuber,
		``Long Short-Term Memory,''
		Neural Computation,
		vol.\ 9,
		no.\ 8,
		pp.\ 1735--1780,
		1997.
		
		\bibitem{cho2014gru}
		K.\ Cho, B.\ van Merri\"{e}nboer, C.\ Gulcehre, D.\ Bahdanau, F.\ Bougares, H.\ Schwenk and Y.\ Bengio,
		``Learning Phrase Representations Using RNN Encoder--Decoder for Statistical Machine Translation,''
		Proceedings of the 2014 Conference on Empirical Methods in Natural Language Processing (EMNLP),
		Doha, Qatar,
		2014,
		pp.\ 1724--1734.
		
		\bibitem{grurnn2018}
		T.A.\ Tang, L.\ Mhamdi, D.\ McLernon, S.A.R.\ Zaidi and M.\ Ghogho,
		``Deep Recurrent Neural Network for Intrusion Detection in SDN-Based Networks,''
		IEEE Conference on Network Softwarization (NetSoft),
		Montreal, QC, Canada,
		2018,
		pp.\ 202--206.
		
		\bibitem{cnnlstm2024}
		P.\ Kaur and H.\ Kumar and M.\ Singh,
		``Enhancing Intrusion Detection Through a Hybrid Deep Learning Pipeline on CIC-IDS2017,''
		Journal of Cloud Computing,
		vol.\ 13,
		Article 94,
		2024.
		
		\bibitem{cnnbigru2022}
		Z.\ Wang and Y.\ Liu and D.\ He and S.\ Chan,
		``Network Intrusion Detection With CNN and BiGRU Under Hybrid Class-Balancing,''
		Security and Communication Networks,
		vol.\ 2022,
		Article 4778735,
		2022.
		
		\bibitem{rtids2022}
		Z.\ Wu and H.\ Zhang and P.\ Wang and Z.\ Sun,
		``RTIDS: A Robust Transformer-Based Intrusion Detection System,''
		IEEE Access,
		vol.\ 10,
		pp.\ 65126--65140,
		2022.
		
		\bibitem{flowformer2024}
		L.D.\ Manocchio and S.\ Layeghy and W.W.\ Lo and G.K.\ Kulatilleke and M.\ Sarhan and M.\ Portmann,
		``FlowTransformer: A Transformer Framework for Flow-Based Network Intrusion Detection Systems,''
		Expert Systems with Applications,
		vol.\ 241,
		Article 122666,
		2024.
		
		\bibitem{long2024transformer}
		Z.\ Long and H.\ Yan and G.\ Shen and X.\ Zhang and H.\ He and L.\ Cheng,
		``A Transformer-Based Network Intrusion Detection Approach for Cloud Security,''
		Journal of Cloud Computing,
		vol.\ 13,
		Article 5,
		2024.
		
		\bibitem{cicids2017}
		I.\ Sharafaldin, A.H.\ Lashkari and A.A.\ Ghorbani,
		``Toward Generating a New Intrusion Detection Dataset and Intrusion Traffic Characterization,''
		2018 4th International Conference on Information Systems Security and Privacy (ICISSP),
		Funchal, Portugal,
		2018,
		pp.\ 108--116.
		
		\bibitem{mlp2021}
		B.\ Mohammed and E.K.\ Gbashi,
		``Intrusion Detection System for NSL-KDD Dataset Based on Deep Learning and Recursive Feature Elimination,''
		Engineering and Technology Journal,
		vol.\ 39,
		no.\ 7,
		pp.\ 1069--1079,
		2021.
		
		\bibitem{breiman2001}
		L.\ Breiman,
		``Random Forests,''
		Machine Learning,
		vol.\ 45,
		no.\ 1,
		pp.\ 5--32,
		2001.
		
		\bibitem{ring2019survey}
		M.\ Ring and S.\ Wunderlich and D.\ Scheuring and D.\ Landes and A.\ Hotho,
		``A Survey of Network-Based Intrusion Detection Data Sets,''
		Computers \& Security,
		vol.\ 86,
		pp.\ 147--167,
		2019.
		
		\bibitem{thakkar2020}
		A.\ Thakkar and R.\ Lohiya,
		``A Review of the Advancement in Intrusion Detection Datasets,''
		Procedia Computer Science,
		vol.\ 167,
		pp.\ 636--645,
		2020.
		
		\bibitem{kurniabudi2020}
		Kurniabudi and D.\ Stiawan and Darmawijoyo and M.Y.\ Idris and A.M.\ Bamhdi and R.\ Budiarto,
		``CICIDS-2017 Dataset Feature Analysis With Information Gain for Anomaly Detection,''
		IEEE Access,
		vol.\ 8,
		pp.\ 132911--132921,
		2020.
		
		\bibitem{xgb2023}
		M.\ Azam and T.\ Sharif and H.\ Aljuaid,
		``XGBoost-Based Intrusion Detection System for the CIC-IDS2017 Benchmark,''
		IEEE Access,
		vol.\ 11,
		pp.\ 104037--104052,
		2023.
		
		\bibitem{chen2016}
		T.\ Chen and C.\ Guestrin,
		``XGBoost: A Scalable Tree Boosting System,''
		Proceedings of the 22nd ACM SIGKDD International Conference on Knowledge Discovery and Data Mining,
		San Francisco, CA, USA,
		2016,
		pp.\ 785--794.
		
		\bibitem{svmelm2018}
		I.\ Ahmad and M.\ Basheri and M.J.\ Iqbal and A.\ Rahim,
		``Performance Comparison of Support Vector Machine, Extreme Learning Machine, and Deep Learning Classifiers for Network Intrusion Detection,''
		IEEE Access,
		vol.\ 6,
		pp.\ 33789--33795,
		2018.
		
		\bibitem{cnn2019}
		R.\ Abdulhammed and H.\ Musafer and A.\ Alessa and M.\ Faezipour and A.\ Abuzneid,
		``Features Dimensionality Reduction Approaches for Machine Learning Based Network Intrusion Detection Systems,''
		Electronics,
		vol.\ 8,
		no.\ 3,
		Article 322,
		2019.
		
		\bibitem{jaeger2001}
		H.\ Jaeger,
		``The Echo State Approach to Analysing and Training Recurrent Neural Networks,''
		German National Research Center for Information Technology (GMD),
		Technical Report 148,
		2001.
		
		\bibitem{lukosevicius2012}
		M.\ Luko\v{s}evi\v{c}ius and H.\ Jaeger,
		``Reservoir Computing Approaches to Recurrent Neural Network Training,''
		Computer Science Review,
		vol.\ 3,
		no.\ 3,
		pp.\ 127--149,
		2009.
		
		\bibitem{esn_dos2021}
		K.\ Bekshentayeva and L.\ Trajkovi\'{c},
		``Detection of Denial-of-Service Attacks Using Echo State Networks,''
		IEEE International Conference on Systems, Man, and Cybernetics (SMC),
		Melbourne, Australia,
		2021,
		pp.\ 2085--2090.
		
		\bibitem{mlesn2020}
		S.T.\ Zhang and Z.H.\ Liang and Y.\ Wang and X.\ Zhao,
		``Network Traffic Anomaly Detection Based on ML-ESN for Power Metering Systems,''
		Mathematical Problems in Engineering,
		vol.\ 2020,
		Article 8892237,
		2020.
		
		\bibitem{gallicchio2017}
		C.\ Gallicchio and A.\ Micheli and L.\ Pedrelli,
		``Deep Reservoir Computing: A Critical Experimental Analysis,''
		Neurocomputing,
		vol.\ 268,
		pp.\ 87--99,
		2017.
		
		\bibitem{lopez2018lstmesn}
		E.\ L\'{o}pez, C.\ Valle, H.\ Allende, E.\ Gil and H.\ Madsen,
		``Wind Power Forecasting Based on Echo State Networks and Long Short-Term Memory,''
		Energies,
		vol.\ 11,
		no.\ 3,
		Article 526,
		2018.
		
		\bibitem{cesn2024}
		S.\ Bahmaid and S.A.M.\ Ghaleb,
		``Intrusion Detection System Using Chaotic Walrus Optimization-Based Convolutional Echo State Networks for IoT-Assisted Wireless Sensor Networks,''
		Journal of Wireless Mobile Networks, Ubiquitous Computing, and Dependable Applications,
		vol.\ 15,
		no.\ 3,
		pp.\ 236--252,
		2024.
		
		\bibitem{vaswani2017attention}
		A.\ Vaswani, N.\ Shazeer, N.\ Parmar, J.\ Uszkoreit, L.\ Jones, A.N.\ Gomez, L.\ Kaiser and I.\ Polosukhin,
		``Attention Is All You Need,''
		Advances in Neural Information Processing Systems (NeurIPS),
		Long Beach, CA, USA,
		2017,
		pp.\ 5998--6008.
		
		\bibitem{bilstm_attn2023}
		J.\ Zhang and X.\ Zhang and Z.\ Liu and F.\ Fu and Y.\ Jiao and F.\ Xu,
		``A Network Intrusion Detection Model Based on BiLSTM With Multi-Head Attention Mechanism,''
		Electronics,
		vol.\ 12,
		no.\ 19,
		Article 4170,
		2023.
		
		\bibitem{trbma2025}
		Y.\ Li and X.\ Chen and W.\ Zhang,
		``TRBMA: Temporal ResNet Coupled With Bidirectional Multi-Head Attention for Intrusion Detection on CIC-IDS2017,''
		PLoS ONE,
		vol.\ 20,
		no.\ 2,
		Article e0318453,
		2025.
		
		\bibitem{kheddar2024transformers}
		H.\ Kheddar,
		``Transformers and Large Language Models for Efficient Intrusion Detection Systems: A Comprehensive Survey,''
		arXiv preprint arXiv:2408.07583,
		2024.
		
		\bibitem{oh2025mchensemble}
		S.\ Oh and S.\ Sohn and C.\ Kim and M.\ Park,
		``MCH-Ensemble: Minority Class Highlighting Ensemble Method for Class Imbalance in Network Intrusion Detection,''
		Applied Sciences,
		vol.\ 15,
		no.\ 23,
		Article 12647,
		2025.
		
		\bibitem{hambali2026breaking}
		M.A.\ Hambali and N.Z.\ Bako and M.\ Dalhatu and A.\ Ishaq,
		``Breaking Class Imbalance Barriers in Intrusion Detection Systems: A Clustering-Based Hybrid Framework,''
		Scientific Journal of Computer Science,
		vol.\ 2,
		no.\ 1,
		pp.\ 81--95,
		2026.
		
		\bibitem{scardapane2017}
		S.\ Scardapane and D.\ Wang,
		``Randomness in Neural Networks: An Overview,''
		WIREs Data Mining and Knowledge Discovery,
		vol.\ 7,
		no.\ 2,
		Article e1200,
		2017.
		
		\bibitem{rf2020}
		Kurniabudi and D.\ Stiawan and Darmawijoyo and M.Y.\ Idris and A.M.\ Bamhdi and R.\ Budiarto,
		``CICIDS-2017 Dataset Feature Analysis With Information Gain for Anomaly Detection,''
		IEEE Access,
		vol.\ 8,
		pp.\ 132911--132921,
		2020.
		
		\bibitem{reis2019}
		B.\ Reis and E.\ Maia and I.\ Praça,
		``Selection and Performance Analysis of CICIDS2017 Features Importance,''
		International Symposium on Foundations and Practice of Security (FPS),
		Paris, France,
		2019,
		pp.\ 56--71.
		
		\bibitem{engelen2021}
		G.\ Engelen and V.\ Rimmer and W.\ Joosen,
		``Troubleshooting an Intrusion Detection Dataset: The CICIDS2017 Case Study,''
		IEEE Symposium on Security and Privacy Workshops (SPW),
		San Francisco, CA, USA,
		2021,
		pp.\ 7--12.
		
		\bibitem{lashkari2017}
		A.H.\ Lashkari and G.D.\ Gil and M.S.I.\ Mamun and A.A.\ Ghorbani,
		``Characterization of Tor Traffic Using Time-Based Features,''
		2017 3rd International Conference on Information Systems Security and Privacy (ICISSP),
		Porto, Portugal,
		2017,
		pp.\ 253--262.
		
		\bibitem{alam2026drlesn}
		K.\ Alam and M.H.\ Bhuiyan and D.M.\ Farid,
		``A Deep Reinforcement-Based Echo State Network for Network Intrusion Classification,''
		PLOS ONE,
		vol.\ 21,
		no.\ 4,
		Article e0333038,
		2026.
		
		\bibitem{idrissi2023fedanids}
		M.J.\ Idrissi and H.\ Alami and A.\ El Mahdaouy and A.\ El Mekki and S.\ Oualil and Z.\ Yartaoui and I.\ Berrada,
		``Fed-ANIDS: Federated Learning for Anomaly-Based Network Intrusion Detection Systems,''
		Expert Systems with Applications,
		vol.\ 234,
		Article 121000,
		2023.
		
		\bibitem{feddef2023}
		J.\ Chen and Y.\ Zhao and Q.\ Li and X.\ Feng and K.\ Xu,
		``FedDef: Defense Against Gradient Leakage in Federated Learning-Based Network Intrusion Detection Systems,''
		IEEE Transactions on Information Forensics and Security,
		vol.\ 18,
		pp.\ 4561--4576,
		2023.
		
		\bibitem{hore2024sequential}
		S.\ Hore and J.\ Ghadermazi and A.\ Shah and N.D.\ Bastian,
		``A Sequential Deep Learning Framework for a Robust and Resilient Network Intrusion Detection System,''
		Computers \& Security,
		vol.\ 144,
		Article 103928,
		2024.
		
		\bibitem{hernandezramos2025federated}
		J.L.\ Hernandez-Ramos and G.\ Karopoulos and E.\ Chatzoglou and V.\ Kouliaridis and E.\ Marmol and A.\ Gonzalez-Vidal and G.\ Kambourakis,
		``Intrusion Detection Based on Federated Learning: A Systematic Review,''
		ACM Computing Surveys,
		vol.\ 57,
		no.\ 12,
		Article 309,
		2025.
		
		\bibitem{almansour2025adaptive}
		S.\ Almansour and K.\ Yadav and L.M.\ Alkwai and N.S.\ Alghamdi and W.\ Viriyasitavat and G.\ Dhiman,
		``Adaptive Personalized Federated Learning With Lightweight Depthwise Convolutional Bottleneck Network for Novel Intrusion Detection System in Internet of Vehicles,''
		Scientific Reports,
		vol.\ 15,
		Article 35604,
		2025.
		
		\bibitem{moczkodan2026transformers}
		Z.\ Moczkodan and H.\ Ragab,
		``Do Transformers Actually Help Intrusion Detection? A Temporal Sequence Evaluation on CIC-IDS2017,''
		arXiv preprint arXiv:2606.11098,
		2026.
		
		\bibitem{zhou2020ensemble}
		Y.\ Zhou and G.\ Cheng and S.\ Jiang and M.\ Dai,
		``Building an Efficient Intrusion Detection System Based on Feature Selection and Ensemble Classifier,''
		Computer Networks,
		vol.\ 174,
		Article 107247,
		2020.
		
		\bibitem{gallicchio2020}
		C.\ Gallicchio and A.\ Micheli,
		``Ring Reservoir Neural Networks for Graphs,''
		IEEE International Joint Conference on Neural Networks (IJCNN),
		Glasgow, United Kingdom,
		2020,
		pp.\ 1--7.
		
		\bibitem{liu2020hierarchical}
		C.\ Liu and Y.\ Liu and Y.\ Yan and J.\ Wang,
		``An Intrusion Detection Model With Hierarchical Attention Mechanism,''
		IEEE Access,
		vol.\ 8,
		pp.\ 67542--67554,
		2020.
		
	\end{thebibliography}
\end{document}